% The marks of the series: what a paper is, in the left margin of its first page, and, while the
% series is a draft, the word DRAFT across every page of it.
%
% One file for the three papers, so that a later change to the wording, the colour or the position
% lands in all three at once. It draws, rotated a quarter turn and read from the bottom, the state
% of the paper, the repository, the paper's own folder and the date of the revision it prints:
%
%     DRAFT . evilaliv3/vesuvius-challenge-pipeline . research/00001 . last updated 18 September 2026
%
% Why it exists. These are self published reports whose claims are pinned to a codebase that moves
% under them and to data that changes, and the method they argue for rests on when a thing was
% declared against when it was measured. A paper that hides its own date is inconsistent with the
% way it argues, and a reader who prints the front page should be able to see, without reading a
% line of it, what the thing is, what state it is in and when it was last revised.
%
% Usage, two lines per paper. In the PREAMBLE, after the class and before \begin{document}:
%
%     % The marks of the series: what a paper is, in the left margin of its first page, and, while the
% series is a draft, the word DRAFT across every page of it.
%
% One file for the three papers, so that a later change to the wording, the colour or the position
% lands in all three at once. It draws, rotated a quarter turn and read from the bottom, the state
% of the paper, the repository, the paper's own folder and the date of the revision it prints:
%
%     DRAFT . evilaliv3/vesuvius-challenge-pipeline . research/00001 . last updated 18 September 2026
%
% Why it exists. These are self published reports whose claims are pinned to a codebase that moves
% under them and to data that changes, and the method they argue for rests on when a thing was
% declared against when it was measured. A paper that hides its own date is inconsistent with the
% way it argues, and a reader who prints the front page should be able to see, without reading a
% line of it, what the thing is, what state it is in and when it was last revised.
%
% Usage, two lines per paper. In the PREAMBLE, after the class and before \begin{document}:
%
%     % The marks of the series: what a paper is, in the left margin of its first page, and, while the
% series is a draft, the word DRAFT across every page of it.
%
% One file for the three papers, so that a later change to the wording, the colour or the position
% lands in all three at once. It draws, rotated a quarter turn and read from the bottom, the state
% of the paper, the repository, the paper's own folder and the date of the revision it prints:
%
%     DRAFT . evilaliv3/vesuvius-challenge-pipeline . research/00001 . last updated 18 September 2026
%
% Why it exists. These are self published reports whose claims are pinned to a codebase that moves
% under them and to data that changes, and the method they argue for rests on when a thing was
% declared against when it was measured. A paper that hides its own date is inconsistent with the
% way it argues, and a reader who prints the front page should be able to see, without reading a
% line of it, what the thing is, what state it is in and when it was last revised.
%
% Usage, two lines per paper. In the PREAMBLE, after the class and before \begin{document}:
%
%     % The marks of the series: what a paper is, in the left margin of its first page, and, while the
% series is a draft, the word DRAFT across every page of it.
%
% One file for the three papers, so that a later change to the wording, the colour or the position
% lands in all three at once. It draws, rotated a quarter turn and read from the bottom, the state
% of the paper, the repository, the paper's own folder and the date of the revision it prints:
%
%     DRAFT . evilaliv3/vesuvius-challenge-pipeline . research/00001 . last updated 18 September 2026
%
% Why it exists. These are self published reports whose claims are pinned to a codebase that moves
% under them and to data that changes, and the method they argue for rests on when a thing was
% declared against when it was measured. A paper that hides its own date is inconsistent with the
% way it argues, and a reader who prints the front page should be able to see, without reading a
% line of it, what the thing is, what state it is in and when it was last revised.
%
% Usage, two lines per paper. In the PREAMBLE, after the class and before \begin{document}:
%
%     \input{sidestripe}
%
% and in the BODY, immediately after \begin{document}, with the paper's own folder number:
%
%     \sidestripe{00001}
%
% The second line has to be in the body because the stripe is attached to the next page shipped
% out, which is the first page; the first line has to be in the preamble because it loads packages.
% The build of each paper fails if it carries neither.
%
% The date is \BuildDateLong, written by the build from SOURCE_DATE_EPOCH, the same value pdfTeX
% stamps into /CreationDate. It is never typed, so the front page and the metadata cannot come to
% say two different days. It is the date of the revision and not a reading of the clock, which is
% why the stripe labels it "last updated": a bare date beside a title is read as the day of
% writing, of submission or of the build, and a reader has no way to tell which was meant.

% THE DRAFT SWITCH OF THE SERIES. This one line is the whole control, for all three papers.
%
%     \SeriesDrafttrue    the papers are drafts and carry the mark
%     \SeriesDraftfalse   the papers are final and carry nothing
%
% Nothing else moves when it is thrown. The three papers reach this file through
% \input{sidestripe}, and the build of each one refuses a paper that does not, so none of them can
% miss the switch and no paper of the series can be marked while another is not.
%
% Off, this file draws no watermark and loads no package it did not load before, and each paper
% rebuilds to the bytes it had before the mark existed. That was checked on 2026-09-18 by building
% all three with the switch off and comparing the sha256 with the published files, not assumed.
\newif\ifSeriesDraft
\SeriesDraftfalse

\usepackage{eso-pic}
\usepackage{xcolor}

% The build writes build-date.tex and every paper of the series inputs it before this file. If it
% is missing, the date is not declared anywhere, and the build stops here rather than reaching a
% page that quietly carries no date or, worse, one taken from the clock.
\ifdefined\BuildDateLong\else
  \PackageError{sidestripe}{The command BuildDateLong is not defined: input build-date first}
    {The build writes build-date.tex from SOURCE_DATE_EPOCH. Run the paper's build script.}
\fi

\newlength{\SideStripeX}\setlength{\SideStripeX}{0.42in}
\newlength{\SideStripeY}\setlength{\SideStripeY}{1.05in}
% The separator is a middle dot: the house rule forbids a dash as punctuation, and a comma would
% read as a list rather than as the three fields of an identifier.
\newcommand{\SideStripeSep}{\,\textperiodcentered\,}
\newcommand{\SideStripeRepo}{evilaliv3/vesuvius-challenge-pipeline}

% The watermark, drawn only while the switch is on.
%
% Why a mark on every page and not a banner on the first. A reader almost never meets these papers
% whole: a page is printed, a screenshot is pasted into a message, a section is forwarded. A mark
% that lives on the first page alone is gone at that moment, and the page that travels is the one
% that will be quoted back. A running head would reach every page too, but it belongs to the
% galley: in this two column class it would move the text block, change where the columns break
% and cost a page, and it would have to be removed from the galley again to get the final paper
% back. A shipout picture costs nothing of the kind.
%
% Why the FOREGROUND and not the background. It was drawn in the background until 2026-09-18, and
% there it was invisible wherever a figure sat, because every figure in these papers fills its own
% box with opaque white before it draws anything. A figure is exactly the part of a page somebody
% screenshots, so the mark was missing from the part most likely to travel alone. On top of
% everything it reaches the whole page, figures included.
%
% What keeps it from spoiling the page. It is drawn at \SeriesDraftOpacity of full opacity, so what
% lands on the paper is a tint and not a shape: over white it is the lightest grey a PDF can carry,
% and over black text it changes nothing at all, because a tint of black over black is black. The
% text underneath is therefore not dimmed in any way a reader can measure. Nothing is added to the
% galley either, so the page count, the line breaks and the column breaks are exactly those of the
% final paper, and the gate on overfull boxes sees nothing new.
%
% The colour is black and not a grey, because at this opacity the colour is the only thing left to
% spend: the tint over white is 1 minus the opacity of the way from white to the ink, so black is
% the darkest mark two per cent can buy. A light grey at two per cent would be indistinguishable
% from the paper.
%
% \resizebox scales a box set at the normal size instead of asking for a font at 80pt, which the
% Computer Modern shapes of this class do not have: LaTeX would substitute the nearest design size
% and print a mark several times smaller than the one asked for, with a font warning the gates do
% not read. Scaling the box cannot fail that way, and it is a linear transform in the PDF, so the
% result is the same bytes on every build.
\ifSeriesDraft
  % transparent.sty writes the constant alpha into an ExtGState, which is what pdfTeX needs to put
  % ink on top of ink without hiding it. It is loaded only here, so a paper built with the switch
  % off carries no transparency machinery and no extra resource at all.
  \usepackage{transparent}
  \newcommand{\SeriesDraftWord}{DRAFT}
  \newcommand{\SeriesDraftOpacity}{0.10}
  \AddToShipoutPictureFG{%
    \AtPageCenter{%
      \makebox(0,0){%
        \rotatebox{45}{%
          \transparent{\SeriesDraftOpacity}%
          \resizebox{0.68\paperwidth}{!}{%
            \color{black}\sffamily\bfseries\SeriesDraftWord}}}}}
\fi

% THE CASE OF THE SERIES: the headings and the table captions read as the rest of the article.
%
% IEEEtran sets two things in small capitals, and a reader sees them as ALL CAPITALS shouted in the
% middle of a page whose every other line is ordinary sentence case:
%
%   the table captions, IEEEtran.cls line 2778, which puts the caption text through \scshape;
%   the section headings, IEEEtran.cls line 5473, whose heading font ends in \scshape.
%
% Those are the two the responsible person asked for on 2026-09-18, and they are the only two that
% are live in this class option. The other \scshape sites in the class, lines 2737, 5468, 5501 and
% 5790, belong to the conference and computer society options, which these papers do not use, and
% the figure captions never had it: figure captions were already sentence case while table captions
% were not, so this change also makes the two agree, which they did not before.
%
% What is deliberately NOT changed: the drop cap at the start of each section, IEEEtran.cls lines
% 5835 and 5841, which sets the large initial and the rest of its first word through
% \MakeUppercase. A large initial followed by a word in capitals is a typographic convention and
% not shouting, and it is a bigger aesthetic move than was asked for. One line in this file takes
% it off if it is ever wanted: \renewcommand{\IEEEPARstartWORDCAPSTYLE}{\relax}.
%
% The class file is NOT edited. It is a third party class, vendored in three trees, and a modified
% copy under the same name is what makes a build unreproducible for somebody else. etoolbox patches
% the two macros from the preamble instead, so one change here reaches all three papers exactly as
% the draft switch does.
%
% Both patches fail the build if they do not apply. That matters more than it looks: if the class
% is ever updated and the macro no longer matches, a silent failure would bring the capitals back
% and nobody would know which build changed them.
\usepackage{etoolbox}
\makeatletter
\patchcmd{\@makecaption}{\scshape}{}{}{%
  \PackageError{sidestripe}{Cannot take the small capitals out of the table captions}
    {IEEEtran's \string\@makecaption no longer contains \string\scshape where it did at line
     2778. Read the class and fix the patch rather than removing it.}}
\patchcmd{\section}{\scshape}{}{}{%
  \PackageError{sidestripe}{Cannot take the small capitals out of the section headings}
    {IEEEtran's \string\section no longer contains \string\scshape where it did at line 5473.
     Read the class and fix the patch rather than removing it.}}
\makeatother

% Mechanics of the stripe. eso-pic's \AddToShipoutPictureBG* draws in the background layer of the
% current page only, which is what keeps the body untouched: the text block is not moved by a
% single point, because nothing is added to the galley. \rotatebox comes from graphicx, which the
% papers already load for their figures, and the grey is xcolor. tikz with `remember picture,
% overlay' would do the same thing at the price of a second compilation pass and a much larger
% dependency, and a raw \rotatebox in a zero size box inside the body would have to be placed by
% hand relative to the text and would move with it.
%
% Geometry. The stripe sits \SideStripeX from the left edge of the paper and starts \SideStripeY
% from the bottom, reading upwards. Those are far enough inside the paper that a printer's trim
% does not eat them and far enough from the text block that they cannot collide with it; both are
% named here so that the three papers cannot drift apart.
% If the paper that inputs this file has no published-date.tex, say so in words rather than
% letting TeX report an undefined control sequence from inside a rotated box. The three papers of
% paper/00001 to 00003 are in exactly that state on 2026-09-22 and are not rebuilt today; whoever
% builds one next is told here what to add and why, instead of debugging it.
\ifdefined\PublishedStripe\else
  \errmessage{sidestripe.tex: this paper inputs the stripe but defines no \string\PublishedStripe.
  Add a published-date.tex beside the paper and \string\input{published-date}. It holds the day
  the folder entered the PUBLIC repository, written by hand and never from the clock; a folder not
  yet published defines \string\PublishedStripe as: not yet published. See REVISION.md,
  2026-09-22}
\fi

\newcommand{\sidestripe}[1]{%
  \AddToShipoutPictureBG*{%
    \AtPageLowerLeft{%
      \put(\LenToUnit{\SideStripeX},\LenToUnit{\SideStripeY}){%
        \rotatebox{90}{%
          \color[gray]{0.55}%
          \fontsize{7}{8}\selectfont\ttfamily
          \ifSeriesDraft\SeriesDraftWord\SideStripeSep\fi%
          \SideStripeRepo\SideStripeSep results/#1\SideStripeSep
          \PublishedStripe\SideStripeSep last updated \BuildDateLong}}}}}

% CHANGE OF 2026-09-22, on the owner's word, to the one file all the papers of the series share.
%
% The stripe now carries TWO dates and names the folder as results/ rather than research/:
%
%   DRAFT . evilaliv3/vesuvius-challenge-pipeline . results/<folder> . <published> . last updated <date>
%
% \PublishedStripe comes from a published-date.tex beside the paper and is the day the folder
% entered the PUBLIC repository, written by hand once and never taken from the clock; a folder
% that is not published yet defines it as «not yet published». \BuildDateLong is the other
% quantity and keeps coming from SOURCE_DATE_EPOCH. A build refuses a publication date that is
% after the build date.
%
% EVERY PAPER THAT INPUTS THIS FILE MUST NOW ALSO INPUT A published-date.tex, or \PublishedStripe
% is undefined and the build stops. The three published papers of paper/00001 to 00003 are NOT
% rebuilt today: they are recompiled with their true publication dates at their next revision,
% which is the owner's decision of 2026-09-22 and is written in REVISION.md with its date.

%
% and in the BODY, immediately after \begin{document}, with the paper's own folder number:
%
%     \sidestripe{00001}
%
% The second line has to be in the body because the stripe is attached to the next page shipped
% out, which is the first page; the first line has to be in the preamble because it loads packages.
% The build of each paper fails if it carries neither.
%
% The date is \BuildDateLong, written by the build from SOURCE_DATE_EPOCH, the same value pdfTeX
% stamps into /CreationDate. It is never typed, so the front page and the metadata cannot come to
% say two different days. It is the date of the revision and not a reading of the clock, which is
% why the stripe labels it "last updated": a bare date beside a title is read as the day of
% writing, of submission or of the build, and a reader has no way to tell which was meant.

% THE DRAFT SWITCH OF THE SERIES. This one line is the whole control, for all three papers.
%
%     \SeriesDrafttrue    the papers are drafts and carry the mark
%     \SeriesDraftfalse   the papers are final and carry nothing
%
% Nothing else moves when it is thrown. The three papers reach this file through
% % The marks of the series: what a paper is, in the left margin of its first page, and, while the
% series is a draft, the word DRAFT across every page of it.
%
% One file for the three papers, so that a later change to the wording, the colour or the position
% lands in all three at once. It draws, rotated a quarter turn and read from the bottom, the state
% of the paper, the repository, the paper's own folder and the date of the revision it prints:
%
%     DRAFT . evilaliv3/vesuvius-challenge-pipeline . research/00001 . last updated 18 September 2026
%
% Why it exists. These are self published reports whose claims are pinned to a codebase that moves
% under them and to data that changes, and the method they argue for rests on when a thing was
% declared against when it was measured. A paper that hides its own date is inconsistent with the
% way it argues, and a reader who prints the front page should be able to see, without reading a
% line of it, what the thing is, what state it is in and when it was last revised.
%
% Usage, two lines per paper. In the PREAMBLE, after the class and before \begin{document}:
%
%     \input{sidestripe}
%
% and in the BODY, immediately after \begin{document}, with the paper's own folder number:
%
%     \sidestripe{00001}
%
% The second line has to be in the body because the stripe is attached to the next page shipped
% out, which is the first page; the first line has to be in the preamble because it loads packages.
% The build of each paper fails if it carries neither.
%
% The date is \BuildDateLong, written by the build from SOURCE_DATE_EPOCH, the same value pdfTeX
% stamps into /CreationDate. It is never typed, so the front page and the metadata cannot come to
% say two different days. It is the date of the revision and not a reading of the clock, which is
% why the stripe labels it "last updated": a bare date beside a title is read as the day of
% writing, of submission or of the build, and a reader has no way to tell which was meant.

% THE DRAFT SWITCH OF THE SERIES. This one line is the whole control, for all three papers.
%
%     \SeriesDrafttrue    the papers are drafts and carry the mark
%     \SeriesDraftfalse   the papers are final and carry nothing
%
% Nothing else moves when it is thrown. The three papers reach this file through
% \input{sidestripe}, and the build of each one refuses a paper that does not, so none of them can
% miss the switch and no paper of the series can be marked while another is not.
%
% Off, this file draws no watermark and loads no package it did not load before, and each paper
% rebuilds to the bytes it had before the mark existed. That was checked on 2026-09-18 by building
% all three with the switch off and comparing the sha256 with the published files, not assumed.
\newif\ifSeriesDraft
\SeriesDraftfalse

\usepackage{eso-pic}
\usepackage{xcolor}

% The build writes build-date.tex and every paper of the series inputs it before this file. If it
% is missing, the date is not declared anywhere, and the build stops here rather than reaching a
% page that quietly carries no date or, worse, one taken from the clock.
\ifdefined\BuildDateLong\else
  \PackageError{sidestripe}{The command BuildDateLong is not defined: input build-date first}
    {The build writes build-date.tex from SOURCE_DATE_EPOCH. Run the paper's build script.}
\fi

\newlength{\SideStripeX}\setlength{\SideStripeX}{0.42in}
\newlength{\SideStripeY}\setlength{\SideStripeY}{1.05in}
% The separator is a middle dot: the house rule forbids a dash as punctuation, and a comma would
% read as a list rather than as the three fields of an identifier.
\newcommand{\SideStripeSep}{\,\textperiodcentered\,}
\newcommand{\SideStripeRepo}{evilaliv3/vesuvius-challenge-pipeline}

% The watermark, drawn only while the switch is on.
%
% Why a mark on every page and not a banner on the first. A reader almost never meets these papers
% whole: a page is printed, a screenshot is pasted into a message, a section is forwarded. A mark
% that lives on the first page alone is gone at that moment, and the page that travels is the one
% that will be quoted back. A running head would reach every page too, but it belongs to the
% galley: in this two column class it would move the text block, change where the columns break
% and cost a page, and it would have to be removed from the galley again to get the final paper
% back. A shipout picture costs nothing of the kind.
%
% Why the FOREGROUND and not the background. It was drawn in the background until 2026-09-18, and
% there it was invisible wherever a figure sat, because every figure in these papers fills its own
% box with opaque white before it draws anything. A figure is exactly the part of a page somebody
% screenshots, so the mark was missing from the part most likely to travel alone. On top of
% everything it reaches the whole page, figures included.
%
% What keeps it from spoiling the page. It is drawn at \SeriesDraftOpacity of full opacity, so what
% lands on the paper is a tint and not a shape: over white it is the lightest grey a PDF can carry,
% and over black text it changes nothing at all, because a tint of black over black is black. The
% text underneath is therefore not dimmed in any way a reader can measure. Nothing is added to the
% galley either, so the page count, the line breaks and the column breaks are exactly those of the
% final paper, and the gate on overfull boxes sees nothing new.
%
% The colour is black and not a grey, because at this opacity the colour is the only thing left to
% spend: the tint over white is 1 minus the opacity of the way from white to the ink, so black is
% the darkest mark two per cent can buy. A light grey at two per cent would be indistinguishable
% from the paper.
%
% \resizebox scales a box set at the normal size instead of asking for a font at 80pt, which the
% Computer Modern shapes of this class do not have: LaTeX would substitute the nearest design size
% and print a mark several times smaller than the one asked for, with a font warning the gates do
% not read. Scaling the box cannot fail that way, and it is a linear transform in the PDF, so the
% result is the same bytes on every build.
\ifSeriesDraft
  % transparent.sty writes the constant alpha into an ExtGState, which is what pdfTeX needs to put
  % ink on top of ink without hiding it. It is loaded only here, so a paper built with the switch
  % off carries no transparency machinery and no extra resource at all.
  \usepackage{transparent}
  \newcommand{\SeriesDraftWord}{DRAFT}
  \newcommand{\SeriesDraftOpacity}{0.10}
  \AddToShipoutPictureFG{%
    \AtPageCenter{%
      \makebox(0,0){%
        \rotatebox{45}{%
          \transparent{\SeriesDraftOpacity}%
          \resizebox{0.68\paperwidth}{!}{%
            \color{black}\sffamily\bfseries\SeriesDraftWord}}}}}
\fi

% THE CASE OF THE SERIES: the headings and the table captions read as the rest of the article.
%
% IEEEtran sets two things in small capitals, and a reader sees them as ALL CAPITALS shouted in the
% middle of a page whose every other line is ordinary sentence case:
%
%   the table captions, IEEEtran.cls line 2778, which puts the caption text through \scshape;
%   the section headings, IEEEtran.cls line 5473, whose heading font ends in \scshape.
%
% Those are the two the responsible person asked for on 2026-09-18, and they are the only two that
% are live in this class option. The other \scshape sites in the class, lines 2737, 5468, 5501 and
% 5790, belong to the conference and computer society options, which these papers do not use, and
% the figure captions never had it: figure captions were already sentence case while table captions
% were not, so this change also makes the two agree, which they did not before.
%
% What is deliberately NOT changed: the drop cap at the start of each section, IEEEtran.cls lines
% 5835 and 5841, which sets the large initial and the rest of its first word through
% \MakeUppercase. A large initial followed by a word in capitals is a typographic convention and
% not shouting, and it is a bigger aesthetic move than was asked for. One line in this file takes
% it off if it is ever wanted: \renewcommand{\IEEEPARstartWORDCAPSTYLE}{\relax}.
%
% The class file is NOT edited. It is a third party class, vendored in three trees, and a modified
% copy under the same name is what makes a build unreproducible for somebody else. etoolbox patches
% the two macros from the preamble instead, so one change here reaches all three papers exactly as
% the draft switch does.
%
% Both patches fail the build if they do not apply. That matters more than it looks: if the class
% is ever updated and the macro no longer matches, a silent failure would bring the capitals back
% and nobody would know which build changed them.
\usepackage{etoolbox}
\makeatletter
\patchcmd{\@makecaption}{\scshape}{}{}{%
  \PackageError{sidestripe}{Cannot take the small capitals out of the table captions}
    {IEEEtran's \string\@makecaption no longer contains \string\scshape where it did at line
     2778. Read the class and fix the patch rather than removing it.}}
\patchcmd{\section}{\scshape}{}{}{%
  \PackageError{sidestripe}{Cannot take the small capitals out of the section headings}
    {IEEEtran's \string\section no longer contains \string\scshape where it did at line 5473.
     Read the class and fix the patch rather than removing it.}}
\makeatother

% Mechanics of the stripe. eso-pic's \AddToShipoutPictureBG* draws in the background layer of the
% current page only, which is what keeps the body untouched: the text block is not moved by a
% single point, because nothing is added to the galley. \rotatebox comes from graphicx, which the
% papers already load for their figures, and the grey is xcolor. tikz with `remember picture,
% overlay' would do the same thing at the price of a second compilation pass and a much larger
% dependency, and a raw \rotatebox in a zero size box inside the body would have to be placed by
% hand relative to the text and would move with it.
%
% Geometry. The stripe sits \SideStripeX from the left edge of the paper and starts \SideStripeY
% from the bottom, reading upwards. Those are far enough inside the paper that a printer's trim
% does not eat them and far enough from the text block that they cannot collide with it; both are
% named here so that the three papers cannot drift apart.
% If the paper that inputs this file has no published-date.tex, say so in words rather than
% letting TeX report an undefined control sequence from inside a rotated box. The three papers of
% paper/00001 to 00003 are in exactly that state on 2026-09-22 and are not rebuilt today; whoever
% builds one next is told here what to add and why, instead of debugging it.
\ifdefined\PublishedStripe\else
  \errmessage{sidestripe.tex: this paper inputs the stripe but defines no \string\PublishedStripe.
  Add a published-date.tex beside the paper and \string\input{published-date}. It holds the day
  the folder entered the PUBLIC repository, written by hand and never from the clock; a folder not
  yet published defines \string\PublishedStripe as: not yet published. See REVISION.md,
  2026-09-22}
\fi

\newcommand{\sidestripe}[1]{%
  \AddToShipoutPictureBG*{%
    \AtPageLowerLeft{%
      \put(\LenToUnit{\SideStripeX},\LenToUnit{\SideStripeY}){%
        \rotatebox{90}{%
          \color[gray]{0.55}%
          \fontsize{7}{8}\selectfont\ttfamily
          \ifSeriesDraft\SeriesDraftWord\SideStripeSep\fi%
          \SideStripeRepo\SideStripeSep results/#1\SideStripeSep
          \PublishedStripe\SideStripeSep last updated \BuildDateLong}}}}}

% CHANGE OF 2026-09-22, on the owner's word, to the one file all the papers of the series share.
%
% The stripe now carries TWO dates and names the folder as results/ rather than research/:
%
%   DRAFT . evilaliv3/vesuvius-challenge-pipeline . results/<folder> . <published> . last updated <date>
%
% \PublishedStripe comes from a published-date.tex beside the paper and is the day the folder
% entered the PUBLIC repository, written by hand once and never taken from the clock; a folder
% that is not published yet defines it as «not yet published». \BuildDateLong is the other
% quantity and keeps coming from SOURCE_DATE_EPOCH. A build refuses a publication date that is
% after the build date.
%
% EVERY PAPER THAT INPUTS THIS FILE MUST NOW ALSO INPUT A published-date.tex, or \PublishedStripe
% is undefined and the build stops. The three published papers of paper/00001 to 00003 are NOT
% rebuilt today: they are recompiled with their true publication dates at their next revision,
% which is the owner's decision of 2026-09-22 and is written in REVISION.md with its date.
, and the build of each one refuses a paper that does not, so none of them can
% miss the switch and no paper of the series can be marked while another is not.
%
% Off, this file draws no watermark and loads no package it did not load before, and each paper
% rebuilds to the bytes it had before the mark existed. That was checked on 2026-09-18 by building
% all three with the switch off and comparing the sha256 with the published files, not assumed.
\newif\ifSeriesDraft
\SeriesDraftfalse

\usepackage{eso-pic}
\usepackage{xcolor}

% The build writes build-date.tex and every paper of the series inputs it before this file. If it
% is missing, the date is not declared anywhere, and the build stops here rather than reaching a
% page that quietly carries no date or, worse, one taken from the clock.
\ifdefined\BuildDateLong\else
  \PackageError{sidestripe}{The command BuildDateLong is not defined: input build-date first}
    {The build writes build-date.tex from SOURCE_DATE_EPOCH. Run the paper's build script.}
\fi

\newlength{\SideStripeX}\setlength{\SideStripeX}{0.42in}
\newlength{\SideStripeY}\setlength{\SideStripeY}{1.05in}
% The separator is a middle dot: the house rule forbids a dash as punctuation, and a comma would
% read as a list rather than as the three fields of an identifier.
\newcommand{\SideStripeSep}{\,\textperiodcentered\,}
\newcommand{\SideStripeRepo}{evilaliv3/vesuvius-challenge-pipeline}

% The watermark, drawn only while the switch is on.
%
% Why a mark on every page and not a banner on the first. A reader almost never meets these papers
% whole: a page is printed, a screenshot is pasted into a message, a section is forwarded. A mark
% that lives on the first page alone is gone at that moment, and the page that travels is the one
% that will be quoted back. A running head would reach every page too, but it belongs to the
% galley: in this two column class it would move the text block, change where the columns break
% and cost a page, and it would have to be removed from the galley again to get the final paper
% back. A shipout picture costs nothing of the kind.
%
% Why the FOREGROUND and not the background. It was drawn in the background until 2026-09-18, and
% there it was invisible wherever a figure sat, because every figure in these papers fills its own
% box with opaque white before it draws anything. A figure is exactly the part of a page somebody
% screenshots, so the mark was missing from the part most likely to travel alone. On top of
% everything it reaches the whole page, figures included.
%
% What keeps it from spoiling the page. It is drawn at \SeriesDraftOpacity of full opacity, so what
% lands on the paper is a tint and not a shape: over white it is the lightest grey a PDF can carry,
% and over black text it changes nothing at all, because a tint of black over black is black. The
% text underneath is therefore not dimmed in any way a reader can measure. Nothing is added to the
% galley either, so the page count, the line breaks and the column breaks are exactly those of the
% final paper, and the gate on overfull boxes sees nothing new.
%
% The colour is black and not a grey, because at this opacity the colour is the only thing left to
% spend: the tint over white is 1 minus the opacity of the way from white to the ink, so black is
% the darkest mark two per cent can buy. A light grey at two per cent would be indistinguishable
% from the paper.
%
% \resizebox scales a box set at the normal size instead of asking for a font at 80pt, which the
% Computer Modern shapes of this class do not have: LaTeX would substitute the nearest design size
% and print a mark several times smaller than the one asked for, with a font warning the gates do
% not read. Scaling the box cannot fail that way, and it is a linear transform in the PDF, so the
% result is the same bytes on every build.
\ifSeriesDraft
  % transparent.sty writes the constant alpha into an ExtGState, which is what pdfTeX needs to put
  % ink on top of ink without hiding it. It is loaded only here, so a paper built with the switch
  % off carries no transparency machinery and no extra resource at all.
  \usepackage{transparent}
  \newcommand{\SeriesDraftWord}{DRAFT}
  \newcommand{\SeriesDraftOpacity}{0.10}
  \AddToShipoutPictureFG{%
    \AtPageCenter{%
      \makebox(0,0){%
        \rotatebox{45}{%
          \transparent{\SeriesDraftOpacity}%
          \resizebox{0.68\paperwidth}{!}{%
            \color{black}\sffamily\bfseries\SeriesDraftWord}}}}}
\fi

% THE CASE OF THE SERIES: the headings and the table captions read as the rest of the article.
%
% IEEEtran sets two things in small capitals, and a reader sees them as ALL CAPITALS shouted in the
% middle of a page whose every other line is ordinary sentence case:
%
%   the table captions, IEEEtran.cls line 2778, which puts the caption text through \scshape;
%   the section headings, IEEEtran.cls line 5473, whose heading font ends in \scshape.
%
% Those are the two the responsible person asked for on 2026-09-18, and they are the only two that
% are live in this class option. The other \scshape sites in the class, lines 2737, 5468, 5501 and
% 5790, belong to the conference and computer society options, which these papers do not use, and
% the figure captions never had it: figure captions were already sentence case while table captions
% were not, so this change also makes the two agree, which they did not before.
%
% What is deliberately NOT changed: the drop cap at the start of each section, IEEEtran.cls lines
% 5835 and 5841, which sets the large initial and the rest of its first word through
% \MakeUppercase. A large initial followed by a word in capitals is a typographic convention and
% not shouting, and it is a bigger aesthetic move than was asked for. One line in this file takes
% it off if it is ever wanted: \renewcommand{\IEEEPARstartWORDCAPSTYLE}{\relax}.
%
% The class file is NOT edited. It is a third party class, vendored in three trees, and a modified
% copy under the same name is what makes a build unreproducible for somebody else. etoolbox patches
% the two macros from the preamble instead, so one change here reaches all three papers exactly as
% the draft switch does.
%
% Both patches fail the build if they do not apply. That matters more than it looks: if the class
% is ever updated and the macro no longer matches, a silent failure would bring the capitals back
% and nobody would know which build changed them.
\usepackage{etoolbox}
\makeatletter
\patchcmd{\@makecaption}{\scshape}{}{}{%
  \PackageError{sidestripe}{Cannot take the small capitals out of the table captions}
    {IEEEtran's \string\@makecaption no longer contains \string\scshape where it did at line
     2778. Read the class and fix the patch rather than removing it.}}
\patchcmd{\section}{\scshape}{}{}{%
  \PackageError{sidestripe}{Cannot take the small capitals out of the section headings}
    {IEEEtran's \string\section no longer contains \string\scshape where it did at line 5473.
     Read the class and fix the patch rather than removing it.}}
\makeatother

% Mechanics of the stripe. eso-pic's \AddToShipoutPictureBG* draws in the background layer of the
% current page only, which is what keeps the body untouched: the text block is not moved by a
% single point, because nothing is added to the galley. \rotatebox comes from graphicx, which the
% papers already load for their figures, and the grey is xcolor. tikz with `remember picture,
% overlay' would do the same thing at the price of a second compilation pass and a much larger
% dependency, and a raw \rotatebox in a zero size box inside the body would have to be placed by
% hand relative to the text and would move with it.
%
% Geometry. The stripe sits \SideStripeX from the left edge of the paper and starts \SideStripeY
% from the bottom, reading upwards. Those are far enough inside the paper that a printer's trim
% does not eat them and far enough from the text block that they cannot collide with it; both are
% named here so that the three papers cannot drift apart.
% If the paper that inputs this file has no published-date.tex, say so in words rather than
% letting TeX report an undefined control sequence from inside a rotated box. The three papers of
% paper/00001 to 00003 are in exactly that state on 2026-09-22 and are not rebuilt today; whoever
% builds one next is told here what to add and why, instead of debugging it.
\ifdefined\PublishedStripe\else
  \errmessage{sidestripe.tex: this paper inputs the stripe but defines no \string\PublishedStripe.
  Add a published-date.tex beside the paper and \string% The day this folder entered the PUBLIC repository, written here by hand once and never taken
% from the clock: a build that read the date from the machine would move it every time and a
% front page would claim a day it was not published on. The last updated date beside it in the
% stripe is a different quantity and comes from SOURCE_DATE_EPOCH.
%
% This folder is NOT in the public repository yet: on 2026-09-22 the only folder under
% results/ of evilaliv3/vesuvius-challenge-pipeline is e89eb7ad-umbilicus-estimator. So the
% stripe carries «not yet published» rather than a date, and when the folder is published the
% two lines below are replaced with the day it was, in this shape:
%
%   \newcommand{\PublishedDateIso}{2026-09-22}
%   \newcommand{\PublishedStripe}{published 22 September 2026}
%
% build.sh reads \PublishedDateIso and refuses a date in the future.
\newcommand{\PublishedDateIso}{}
\newcommand{\PublishedStripe}{not yet published}
. It holds the day
  the folder entered the PUBLIC repository, written by hand and never from the clock; a folder not
  yet published defines \string\PublishedStripe as: not yet published. See REVISION.md,
  2026-09-22}
\fi

\newcommand{\sidestripe}[1]{%
  \AddToShipoutPictureBG*{%
    \AtPageLowerLeft{%
      \put(\LenToUnit{\SideStripeX},\LenToUnit{\SideStripeY}){%
        \rotatebox{90}{%
          \color[gray]{0.55}%
          \fontsize{7}{8}\selectfont\ttfamily
          \ifSeriesDraft\SeriesDraftWord\SideStripeSep\fi%
          \SideStripeRepo\SideStripeSep results/#1\SideStripeSep
          \PublishedStripe\SideStripeSep last updated \BuildDateLong}}}}}

% CHANGE OF 2026-09-22, on the owner's word, to the one file all the papers of the series share.
%
% The stripe now carries TWO dates and names the folder as results/ rather than research/:
%
%   DRAFT . evilaliv3/vesuvius-challenge-pipeline . results/<folder> . <published> . last updated <date>
%
% \PublishedStripe comes from a published-date.tex beside the paper and is the day the folder
% entered the PUBLIC repository, written by hand once and never taken from the clock; a folder
% that is not published yet defines it as «not yet published». \BuildDateLong is the other
% quantity and keeps coming from SOURCE_DATE_EPOCH. A build refuses a publication date that is
% after the build date.
%
% EVERY PAPER THAT INPUTS THIS FILE MUST NOW ALSO INPUT A published-date.tex, or \PublishedStripe
% is undefined and the build stops. The three published papers of paper/00001 to 00003 are NOT
% rebuilt today: they are recompiled with their true publication dates at their next revision,
% which is the owner's decision of 2026-09-22 and is written in REVISION.md with its date.

%
% and in the BODY, immediately after \begin{document}, with the paper's own folder number:
%
%     \sidestripe{00001}
%
% The second line has to be in the body because the stripe is attached to the next page shipped
% out, which is the first page; the first line has to be in the preamble because it loads packages.
% The build of each paper fails if it carries neither.
%
% The date is \BuildDateLong, written by the build from SOURCE_DATE_EPOCH, the same value pdfTeX
% stamps into /CreationDate. It is never typed, so the front page and the metadata cannot come to
% say two different days. It is the date of the revision and not a reading of the clock, which is
% why the stripe labels it "last updated": a bare date beside a title is read as the day of
% writing, of submission or of the build, and a reader has no way to tell which was meant.

% THE DRAFT SWITCH OF THE SERIES. This one line is the whole control, for all three papers.
%
%     \SeriesDrafttrue    the papers are drafts and carry the mark
%     \SeriesDraftfalse   the papers are final and carry nothing
%
% Nothing else moves when it is thrown. The three papers reach this file through
% % The marks of the series: what a paper is, in the left margin of its first page, and, while the
% series is a draft, the word DRAFT across every page of it.
%
% One file for the three papers, so that a later change to the wording, the colour or the position
% lands in all three at once. It draws, rotated a quarter turn and read from the bottom, the state
% of the paper, the repository, the paper's own folder and the date of the revision it prints:
%
%     DRAFT . evilaliv3/vesuvius-challenge-pipeline . research/00001 . last updated 18 September 2026
%
% Why it exists. These are self published reports whose claims are pinned to a codebase that moves
% under them and to data that changes, and the method they argue for rests on when a thing was
% declared against when it was measured. A paper that hides its own date is inconsistent with the
% way it argues, and a reader who prints the front page should be able to see, without reading a
% line of it, what the thing is, what state it is in and when it was last revised.
%
% Usage, two lines per paper. In the PREAMBLE, after the class and before \begin{document}:
%
%     % The marks of the series: what a paper is, in the left margin of its first page, and, while the
% series is a draft, the word DRAFT across every page of it.
%
% One file for the three papers, so that a later change to the wording, the colour or the position
% lands in all three at once. It draws, rotated a quarter turn and read from the bottom, the state
% of the paper, the repository, the paper's own folder and the date of the revision it prints:
%
%     DRAFT . evilaliv3/vesuvius-challenge-pipeline . research/00001 . last updated 18 September 2026
%
% Why it exists. These are self published reports whose claims are pinned to a codebase that moves
% under them and to data that changes, and the method they argue for rests on when a thing was
% declared against when it was measured. A paper that hides its own date is inconsistent with the
% way it argues, and a reader who prints the front page should be able to see, without reading a
% line of it, what the thing is, what state it is in and when it was last revised.
%
% Usage, two lines per paper. In the PREAMBLE, after the class and before \begin{document}:
%
%     \input{sidestripe}
%
% and in the BODY, immediately after \begin{document}, with the paper's own folder number:
%
%     \sidestripe{00001}
%
% The second line has to be in the body because the stripe is attached to the next page shipped
% out, which is the first page; the first line has to be in the preamble because it loads packages.
% The build of each paper fails if it carries neither.
%
% The date is \BuildDateLong, written by the build from SOURCE_DATE_EPOCH, the same value pdfTeX
% stamps into /CreationDate. It is never typed, so the front page and the metadata cannot come to
% say two different days. It is the date of the revision and not a reading of the clock, which is
% why the stripe labels it "last updated": a bare date beside a title is read as the day of
% writing, of submission or of the build, and a reader has no way to tell which was meant.

% THE DRAFT SWITCH OF THE SERIES. This one line is the whole control, for all three papers.
%
%     \SeriesDrafttrue    the papers are drafts and carry the mark
%     \SeriesDraftfalse   the papers are final and carry nothing
%
% Nothing else moves when it is thrown. The three papers reach this file through
% \input{sidestripe}, and the build of each one refuses a paper that does not, so none of them can
% miss the switch and no paper of the series can be marked while another is not.
%
% Off, this file draws no watermark and loads no package it did not load before, and each paper
% rebuilds to the bytes it had before the mark existed. That was checked on 2026-09-18 by building
% all three with the switch off and comparing the sha256 with the published files, not assumed.
\newif\ifSeriesDraft
\SeriesDraftfalse

\usepackage{eso-pic}
\usepackage{xcolor}

% The build writes build-date.tex and every paper of the series inputs it before this file. If it
% is missing, the date is not declared anywhere, and the build stops here rather than reaching a
% page that quietly carries no date or, worse, one taken from the clock.
\ifdefined\BuildDateLong\else
  \PackageError{sidestripe}{The command BuildDateLong is not defined: input build-date first}
    {The build writes build-date.tex from SOURCE_DATE_EPOCH. Run the paper's build script.}
\fi

\newlength{\SideStripeX}\setlength{\SideStripeX}{0.42in}
\newlength{\SideStripeY}\setlength{\SideStripeY}{1.05in}
% The separator is a middle dot: the house rule forbids a dash as punctuation, and a comma would
% read as a list rather than as the three fields of an identifier.
\newcommand{\SideStripeSep}{\,\textperiodcentered\,}
\newcommand{\SideStripeRepo}{evilaliv3/vesuvius-challenge-pipeline}

% The watermark, drawn only while the switch is on.
%
% Why a mark on every page and not a banner on the first. A reader almost never meets these papers
% whole: a page is printed, a screenshot is pasted into a message, a section is forwarded. A mark
% that lives on the first page alone is gone at that moment, and the page that travels is the one
% that will be quoted back. A running head would reach every page too, but it belongs to the
% galley: in this two column class it would move the text block, change where the columns break
% and cost a page, and it would have to be removed from the galley again to get the final paper
% back. A shipout picture costs nothing of the kind.
%
% Why the FOREGROUND and not the background. It was drawn in the background until 2026-09-18, and
% there it was invisible wherever a figure sat, because every figure in these papers fills its own
% box with opaque white before it draws anything. A figure is exactly the part of a page somebody
% screenshots, so the mark was missing from the part most likely to travel alone. On top of
% everything it reaches the whole page, figures included.
%
% What keeps it from spoiling the page. It is drawn at \SeriesDraftOpacity of full opacity, so what
% lands on the paper is a tint and not a shape: over white it is the lightest grey a PDF can carry,
% and over black text it changes nothing at all, because a tint of black over black is black. The
% text underneath is therefore not dimmed in any way a reader can measure. Nothing is added to the
% galley either, so the page count, the line breaks and the column breaks are exactly those of the
% final paper, and the gate on overfull boxes sees nothing new.
%
% The colour is black and not a grey, because at this opacity the colour is the only thing left to
% spend: the tint over white is 1 minus the opacity of the way from white to the ink, so black is
% the darkest mark two per cent can buy. A light grey at two per cent would be indistinguishable
% from the paper.
%
% \resizebox scales a box set at the normal size instead of asking for a font at 80pt, which the
% Computer Modern shapes of this class do not have: LaTeX would substitute the nearest design size
% and print a mark several times smaller than the one asked for, with a font warning the gates do
% not read. Scaling the box cannot fail that way, and it is a linear transform in the PDF, so the
% result is the same bytes on every build.
\ifSeriesDraft
  % transparent.sty writes the constant alpha into an ExtGState, which is what pdfTeX needs to put
  % ink on top of ink without hiding it. It is loaded only here, so a paper built with the switch
  % off carries no transparency machinery and no extra resource at all.
  \usepackage{transparent}
  \newcommand{\SeriesDraftWord}{DRAFT}
  \newcommand{\SeriesDraftOpacity}{0.10}
  \AddToShipoutPictureFG{%
    \AtPageCenter{%
      \makebox(0,0){%
        \rotatebox{45}{%
          \transparent{\SeriesDraftOpacity}%
          \resizebox{0.68\paperwidth}{!}{%
            \color{black}\sffamily\bfseries\SeriesDraftWord}}}}}
\fi

% THE CASE OF THE SERIES: the headings and the table captions read as the rest of the article.
%
% IEEEtran sets two things in small capitals, and a reader sees them as ALL CAPITALS shouted in the
% middle of a page whose every other line is ordinary sentence case:
%
%   the table captions, IEEEtran.cls line 2778, which puts the caption text through \scshape;
%   the section headings, IEEEtran.cls line 5473, whose heading font ends in \scshape.
%
% Those are the two the responsible person asked for on 2026-09-18, and they are the only two that
% are live in this class option. The other \scshape sites in the class, lines 2737, 5468, 5501 and
% 5790, belong to the conference and computer society options, which these papers do not use, and
% the figure captions never had it: figure captions were already sentence case while table captions
% were not, so this change also makes the two agree, which they did not before.
%
% What is deliberately NOT changed: the drop cap at the start of each section, IEEEtran.cls lines
% 5835 and 5841, which sets the large initial and the rest of its first word through
% \MakeUppercase. A large initial followed by a word in capitals is a typographic convention and
% not shouting, and it is a bigger aesthetic move than was asked for. One line in this file takes
% it off if it is ever wanted: \renewcommand{\IEEEPARstartWORDCAPSTYLE}{\relax}.
%
% The class file is NOT edited. It is a third party class, vendored in three trees, and a modified
% copy under the same name is what makes a build unreproducible for somebody else. etoolbox patches
% the two macros from the preamble instead, so one change here reaches all three papers exactly as
% the draft switch does.
%
% Both patches fail the build if they do not apply. That matters more than it looks: if the class
% is ever updated and the macro no longer matches, a silent failure would bring the capitals back
% and nobody would know which build changed them.
\usepackage{etoolbox}
\makeatletter
\patchcmd{\@makecaption}{\scshape}{}{}{%
  \PackageError{sidestripe}{Cannot take the small capitals out of the table captions}
    {IEEEtran's \string\@makecaption no longer contains \string\scshape where it did at line
     2778. Read the class and fix the patch rather than removing it.}}
\patchcmd{\section}{\scshape}{}{}{%
  \PackageError{sidestripe}{Cannot take the small capitals out of the section headings}
    {IEEEtran's \string\section no longer contains \string\scshape where it did at line 5473.
     Read the class and fix the patch rather than removing it.}}
\makeatother

% Mechanics of the stripe. eso-pic's \AddToShipoutPictureBG* draws in the background layer of the
% current page only, which is what keeps the body untouched: the text block is not moved by a
% single point, because nothing is added to the galley. \rotatebox comes from graphicx, which the
% papers already load for their figures, and the grey is xcolor. tikz with `remember picture,
% overlay' would do the same thing at the price of a second compilation pass and a much larger
% dependency, and a raw \rotatebox in a zero size box inside the body would have to be placed by
% hand relative to the text and would move with it.
%
% Geometry. The stripe sits \SideStripeX from the left edge of the paper and starts \SideStripeY
% from the bottom, reading upwards. Those are far enough inside the paper that a printer's trim
% does not eat them and far enough from the text block that they cannot collide with it; both are
% named here so that the three papers cannot drift apart.
% If the paper that inputs this file has no published-date.tex, say so in words rather than
% letting TeX report an undefined control sequence from inside a rotated box. The three papers of
% paper/00001 to 00003 are in exactly that state on 2026-09-22 and are not rebuilt today; whoever
% builds one next is told here what to add and why, instead of debugging it.
\ifdefined\PublishedStripe\else
  \errmessage{sidestripe.tex: this paper inputs the stripe but defines no \string\PublishedStripe.
  Add a published-date.tex beside the paper and \string\input{published-date}. It holds the day
  the folder entered the PUBLIC repository, written by hand and never from the clock; a folder not
  yet published defines \string\PublishedStripe as: not yet published. See REVISION.md,
  2026-09-22}
\fi

\newcommand{\sidestripe}[1]{%
  \AddToShipoutPictureBG*{%
    \AtPageLowerLeft{%
      \put(\LenToUnit{\SideStripeX},\LenToUnit{\SideStripeY}){%
        \rotatebox{90}{%
          \color[gray]{0.55}%
          \fontsize{7}{8}\selectfont\ttfamily
          \ifSeriesDraft\SeriesDraftWord\SideStripeSep\fi%
          \SideStripeRepo\SideStripeSep results/#1\SideStripeSep
          \PublishedStripe\SideStripeSep last updated \BuildDateLong}}}}}

% CHANGE OF 2026-09-22, on the owner's word, to the one file all the papers of the series share.
%
% The stripe now carries TWO dates and names the folder as results/ rather than research/:
%
%   DRAFT . evilaliv3/vesuvius-challenge-pipeline . results/<folder> . <published> . last updated <date>
%
% \PublishedStripe comes from a published-date.tex beside the paper and is the day the folder
% entered the PUBLIC repository, written by hand once and never taken from the clock; a folder
% that is not published yet defines it as «not yet published». \BuildDateLong is the other
% quantity and keeps coming from SOURCE_DATE_EPOCH. A build refuses a publication date that is
% after the build date.
%
% EVERY PAPER THAT INPUTS THIS FILE MUST NOW ALSO INPUT A published-date.tex, or \PublishedStripe
% is undefined and the build stops. The three published papers of paper/00001 to 00003 are NOT
% rebuilt today: they are recompiled with their true publication dates at their next revision,
% which is the owner's decision of 2026-09-22 and is written in REVISION.md with its date.

%
% and in the BODY, immediately after \begin{document}, with the paper's own folder number:
%
%     \sidestripe{00001}
%
% The second line has to be in the body because the stripe is attached to the next page shipped
% out, which is the first page; the first line has to be in the preamble because it loads packages.
% The build of each paper fails if it carries neither.
%
% The date is \BuildDateLong, written by the build from SOURCE_DATE_EPOCH, the same value pdfTeX
% stamps into /CreationDate. It is never typed, so the front page and the metadata cannot come to
% say two different days. It is the date of the revision and not a reading of the clock, which is
% why the stripe labels it "last updated": a bare date beside a title is read as the day of
% writing, of submission or of the build, and a reader has no way to tell which was meant.

% THE DRAFT SWITCH OF THE SERIES. This one line is the whole control, for all three papers.
%
%     \SeriesDrafttrue    the papers are drafts and carry the mark
%     \SeriesDraftfalse   the papers are final and carry nothing
%
% Nothing else moves when it is thrown. The three papers reach this file through
% % The marks of the series: what a paper is, in the left margin of its first page, and, while the
% series is a draft, the word DRAFT across every page of it.
%
% One file for the three papers, so that a later change to the wording, the colour or the position
% lands in all three at once. It draws, rotated a quarter turn and read from the bottom, the state
% of the paper, the repository, the paper's own folder and the date of the revision it prints:
%
%     DRAFT . evilaliv3/vesuvius-challenge-pipeline . research/00001 . last updated 18 September 2026
%
% Why it exists. These are self published reports whose claims are pinned to a codebase that moves
% under them and to data that changes, and the method they argue for rests on when a thing was
% declared against when it was measured. A paper that hides its own date is inconsistent with the
% way it argues, and a reader who prints the front page should be able to see, without reading a
% line of it, what the thing is, what state it is in and when it was last revised.
%
% Usage, two lines per paper. In the PREAMBLE, after the class and before \begin{document}:
%
%     \input{sidestripe}
%
% and in the BODY, immediately after \begin{document}, with the paper's own folder number:
%
%     \sidestripe{00001}
%
% The second line has to be in the body because the stripe is attached to the next page shipped
% out, which is the first page; the first line has to be in the preamble because it loads packages.
% The build of each paper fails if it carries neither.
%
% The date is \BuildDateLong, written by the build from SOURCE_DATE_EPOCH, the same value pdfTeX
% stamps into /CreationDate. It is never typed, so the front page and the metadata cannot come to
% say two different days. It is the date of the revision and not a reading of the clock, which is
% why the stripe labels it "last updated": a bare date beside a title is read as the day of
% writing, of submission or of the build, and a reader has no way to tell which was meant.

% THE DRAFT SWITCH OF THE SERIES. This one line is the whole control, for all three papers.
%
%     \SeriesDrafttrue    the papers are drafts and carry the mark
%     \SeriesDraftfalse   the papers are final and carry nothing
%
% Nothing else moves when it is thrown. The three papers reach this file through
% \input{sidestripe}, and the build of each one refuses a paper that does not, so none of them can
% miss the switch and no paper of the series can be marked while another is not.
%
% Off, this file draws no watermark and loads no package it did not load before, and each paper
% rebuilds to the bytes it had before the mark existed. That was checked on 2026-09-18 by building
% all three with the switch off and comparing the sha256 with the published files, not assumed.
\newif\ifSeriesDraft
\SeriesDraftfalse

\usepackage{eso-pic}
\usepackage{xcolor}

% The build writes build-date.tex and every paper of the series inputs it before this file. If it
% is missing, the date is not declared anywhere, and the build stops here rather than reaching a
% page that quietly carries no date or, worse, one taken from the clock.
\ifdefined\BuildDateLong\else
  \PackageError{sidestripe}{The command BuildDateLong is not defined: input build-date first}
    {The build writes build-date.tex from SOURCE_DATE_EPOCH. Run the paper's build script.}
\fi

\newlength{\SideStripeX}\setlength{\SideStripeX}{0.42in}
\newlength{\SideStripeY}\setlength{\SideStripeY}{1.05in}
% The separator is a middle dot: the house rule forbids a dash as punctuation, and a comma would
% read as a list rather than as the three fields of an identifier.
\newcommand{\SideStripeSep}{\,\textperiodcentered\,}
\newcommand{\SideStripeRepo}{evilaliv3/vesuvius-challenge-pipeline}

% The watermark, drawn only while the switch is on.
%
% Why a mark on every page and not a banner on the first. A reader almost never meets these papers
% whole: a page is printed, a screenshot is pasted into a message, a section is forwarded. A mark
% that lives on the first page alone is gone at that moment, and the page that travels is the one
% that will be quoted back. A running head would reach every page too, but it belongs to the
% galley: in this two column class it would move the text block, change where the columns break
% and cost a page, and it would have to be removed from the galley again to get the final paper
% back. A shipout picture costs nothing of the kind.
%
% Why the FOREGROUND and not the background. It was drawn in the background until 2026-09-18, and
% there it was invisible wherever a figure sat, because every figure in these papers fills its own
% box with opaque white before it draws anything. A figure is exactly the part of a page somebody
% screenshots, so the mark was missing from the part most likely to travel alone. On top of
% everything it reaches the whole page, figures included.
%
% What keeps it from spoiling the page. It is drawn at \SeriesDraftOpacity of full opacity, so what
% lands on the paper is a tint and not a shape: over white it is the lightest grey a PDF can carry,
% and over black text it changes nothing at all, because a tint of black over black is black. The
% text underneath is therefore not dimmed in any way a reader can measure. Nothing is added to the
% galley either, so the page count, the line breaks and the column breaks are exactly those of the
% final paper, and the gate on overfull boxes sees nothing new.
%
% The colour is black and not a grey, because at this opacity the colour is the only thing left to
% spend: the tint over white is 1 minus the opacity of the way from white to the ink, so black is
% the darkest mark two per cent can buy. A light grey at two per cent would be indistinguishable
% from the paper.
%
% \resizebox scales a box set at the normal size instead of asking for a font at 80pt, which the
% Computer Modern shapes of this class do not have: LaTeX would substitute the nearest design size
% and print a mark several times smaller than the one asked for, with a font warning the gates do
% not read. Scaling the box cannot fail that way, and it is a linear transform in the PDF, so the
% result is the same bytes on every build.
\ifSeriesDraft
  % transparent.sty writes the constant alpha into an ExtGState, which is what pdfTeX needs to put
  % ink on top of ink without hiding it. It is loaded only here, so a paper built with the switch
  % off carries no transparency machinery and no extra resource at all.
  \usepackage{transparent}
  \newcommand{\SeriesDraftWord}{DRAFT}
  \newcommand{\SeriesDraftOpacity}{0.10}
  \AddToShipoutPictureFG{%
    \AtPageCenter{%
      \makebox(0,0){%
        \rotatebox{45}{%
          \transparent{\SeriesDraftOpacity}%
          \resizebox{0.68\paperwidth}{!}{%
            \color{black}\sffamily\bfseries\SeriesDraftWord}}}}}
\fi

% THE CASE OF THE SERIES: the headings and the table captions read as the rest of the article.
%
% IEEEtran sets two things in small capitals, and a reader sees them as ALL CAPITALS shouted in the
% middle of a page whose every other line is ordinary sentence case:
%
%   the table captions, IEEEtran.cls line 2778, which puts the caption text through \scshape;
%   the section headings, IEEEtran.cls line 5473, whose heading font ends in \scshape.
%
% Those are the two the responsible person asked for on 2026-09-18, and they are the only two that
% are live in this class option. The other \scshape sites in the class, lines 2737, 5468, 5501 and
% 5790, belong to the conference and computer society options, which these papers do not use, and
% the figure captions never had it: figure captions were already sentence case while table captions
% were not, so this change also makes the two agree, which they did not before.
%
% What is deliberately NOT changed: the drop cap at the start of each section, IEEEtran.cls lines
% 5835 and 5841, which sets the large initial and the rest of its first word through
% \MakeUppercase. A large initial followed by a word in capitals is a typographic convention and
% not shouting, and it is a bigger aesthetic move than was asked for. One line in this file takes
% it off if it is ever wanted: \renewcommand{\IEEEPARstartWORDCAPSTYLE}{\relax}.
%
% The class file is NOT edited. It is a third party class, vendored in three trees, and a modified
% copy under the same name is what makes a build unreproducible for somebody else. etoolbox patches
% the two macros from the preamble instead, so one change here reaches all three papers exactly as
% the draft switch does.
%
% Both patches fail the build if they do not apply. That matters more than it looks: if the class
% is ever updated and the macro no longer matches, a silent failure would bring the capitals back
% and nobody would know which build changed them.
\usepackage{etoolbox}
\makeatletter
\patchcmd{\@makecaption}{\scshape}{}{}{%
  \PackageError{sidestripe}{Cannot take the small capitals out of the table captions}
    {IEEEtran's \string\@makecaption no longer contains \string\scshape where it did at line
     2778. Read the class and fix the patch rather than removing it.}}
\patchcmd{\section}{\scshape}{}{}{%
  \PackageError{sidestripe}{Cannot take the small capitals out of the section headings}
    {IEEEtran's \string\section no longer contains \string\scshape where it did at line 5473.
     Read the class and fix the patch rather than removing it.}}
\makeatother

% Mechanics of the stripe. eso-pic's \AddToShipoutPictureBG* draws in the background layer of the
% current page only, which is what keeps the body untouched: the text block is not moved by a
% single point, because nothing is added to the galley. \rotatebox comes from graphicx, which the
% papers already load for their figures, and the grey is xcolor. tikz with `remember picture,
% overlay' would do the same thing at the price of a second compilation pass and a much larger
% dependency, and a raw \rotatebox in a zero size box inside the body would have to be placed by
% hand relative to the text and would move with it.
%
% Geometry. The stripe sits \SideStripeX from the left edge of the paper and starts \SideStripeY
% from the bottom, reading upwards. Those are far enough inside the paper that a printer's trim
% does not eat them and far enough from the text block that they cannot collide with it; both are
% named here so that the three papers cannot drift apart.
% If the paper that inputs this file has no published-date.tex, say so in words rather than
% letting TeX report an undefined control sequence from inside a rotated box. The three papers of
% paper/00001 to 00003 are in exactly that state on 2026-09-22 and are not rebuilt today; whoever
% builds one next is told here what to add and why, instead of debugging it.
\ifdefined\PublishedStripe\else
  \errmessage{sidestripe.tex: this paper inputs the stripe but defines no \string\PublishedStripe.
  Add a published-date.tex beside the paper and \string\input{published-date}. It holds the day
  the folder entered the PUBLIC repository, written by hand and never from the clock; a folder not
  yet published defines \string\PublishedStripe as: not yet published. See REVISION.md,
  2026-09-22}
\fi

\newcommand{\sidestripe}[1]{%
  \AddToShipoutPictureBG*{%
    \AtPageLowerLeft{%
      \put(\LenToUnit{\SideStripeX},\LenToUnit{\SideStripeY}){%
        \rotatebox{90}{%
          \color[gray]{0.55}%
          \fontsize{7}{8}\selectfont\ttfamily
          \ifSeriesDraft\SeriesDraftWord\SideStripeSep\fi%
          \SideStripeRepo\SideStripeSep results/#1\SideStripeSep
          \PublishedStripe\SideStripeSep last updated \BuildDateLong}}}}}

% CHANGE OF 2026-09-22, on the owner's word, to the one file all the papers of the series share.
%
% The stripe now carries TWO dates and names the folder as results/ rather than research/:
%
%   DRAFT . evilaliv3/vesuvius-challenge-pipeline . results/<folder> . <published> . last updated <date>
%
% \PublishedStripe comes from a published-date.tex beside the paper and is the day the folder
% entered the PUBLIC repository, written by hand once and never taken from the clock; a folder
% that is not published yet defines it as «not yet published». \BuildDateLong is the other
% quantity and keeps coming from SOURCE_DATE_EPOCH. A build refuses a publication date that is
% after the build date.
%
% EVERY PAPER THAT INPUTS THIS FILE MUST NOW ALSO INPUT A published-date.tex, or \PublishedStripe
% is undefined and the build stops. The three published papers of paper/00001 to 00003 are NOT
% rebuilt today: they are recompiled with their true publication dates at their next revision,
% which is the owner's decision of 2026-09-22 and is written in REVISION.md with its date.
, and the build of each one refuses a paper that does not, so none of them can
% miss the switch and no paper of the series can be marked while another is not.
%
% Off, this file draws no watermark and loads no package it did not load before, and each paper
% rebuilds to the bytes it had before the mark existed. That was checked on 2026-09-18 by building
% all three with the switch off and comparing the sha256 with the published files, not assumed.
\newif\ifSeriesDraft
\SeriesDraftfalse

\usepackage{eso-pic}
\usepackage{xcolor}

% The build writes build-date.tex and every paper of the series inputs it before this file. If it
% is missing, the date is not declared anywhere, and the build stops here rather than reaching a
% page that quietly carries no date or, worse, one taken from the clock.
\ifdefined\BuildDateLong\else
  \PackageError{sidestripe}{The command BuildDateLong is not defined: input build-date first}
    {The build writes build-date.tex from SOURCE_DATE_EPOCH. Run the paper's build script.}
\fi

\newlength{\SideStripeX}\setlength{\SideStripeX}{0.42in}
\newlength{\SideStripeY}\setlength{\SideStripeY}{1.05in}
% The separator is a middle dot: the house rule forbids a dash as punctuation, and a comma would
% read as a list rather than as the three fields of an identifier.
\newcommand{\SideStripeSep}{\,\textperiodcentered\,}
\newcommand{\SideStripeRepo}{evilaliv3/vesuvius-challenge-pipeline}

% The watermark, drawn only while the switch is on.
%
% Why a mark on every page and not a banner on the first. A reader almost never meets these papers
% whole: a page is printed, a screenshot is pasted into a message, a section is forwarded. A mark
% that lives on the first page alone is gone at that moment, and the page that travels is the one
% that will be quoted back. A running head would reach every page too, but it belongs to the
% galley: in this two column class it would move the text block, change where the columns break
% and cost a page, and it would have to be removed from the galley again to get the final paper
% back. A shipout picture costs nothing of the kind.
%
% Why the FOREGROUND and not the background. It was drawn in the background until 2026-09-18, and
% there it was invisible wherever a figure sat, because every figure in these papers fills its own
% box with opaque white before it draws anything. A figure is exactly the part of a page somebody
% screenshots, so the mark was missing from the part most likely to travel alone. On top of
% everything it reaches the whole page, figures included.
%
% What keeps it from spoiling the page. It is drawn at \SeriesDraftOpacity of full opacity, so what
% lands on the paper is a tint and not a shape: over white it is the lightest grey a PDF can carry,
% and over black text it changes nothing at all, because a tint of black over black is black. The
% text underneath is therefore not dimmed in any way a reader can measure. Nothing is added to the
% galley either, so the page count, the line breaks and the column breaks are exactly those of the
% final paper, and the gate on overfull boxes sees nothing new.
%
% The colour is black and not a grey, because at this opacity the colour is the only thing left to
% spend: the tint over white is 1 minus the opacity of the way from white to the ink, so black is
% the darkest mark two per cent can buy. A light grey at two per cent would be indistinguishable
% from the paper.
%
% \resizebox scales a box set at the normal size instead of asking for a font at 80pt, which the
% Computer Modern shapes of this class do not have: LaTeX would substitute the nearest design size
% and print a mark several times smaller than the one asked for, with a font warning the gates do
% not read. Scaling the box cannot fail that way, and it is a linear transform in the PDF, so the
% result is the same bytes on every build.
\ifSeriesDraft
  % transparent.sty writes the constant alpha into an ExtGState, which is what pdfTeX needs to put
  % ink on top of ink without hiding it. It is loaded only here, so a paper built with the switch
  % off carries no transparency machinery and no extra resource at all.
  \usepackage{transparent}
  \newcommand{\SeriesDraftWord}{DRAFT}
  \newcommand{\SeriesDraftOpacity}{0.10}
  \AddToShipoutPictureFG{%
    \AtPageCenter{%
      \makebox(0,0){%
        \rotatebox{45}{%
          \transparent{\SeriesDraftOpacity}%
          \resizebox{0.68\paperwidth}{!}{%
            \color{black}\sffamily\bfseries\SeriesDraftWord}}}}}
\fi

% THE CASE OF THE SERIES: the headings and the table captions read as the rest of the article.
%
% IEEEtran sets two things in small capitals, and a reader sees them as ALL CAPITALS shouted in the
% middle of a page whose every other line is ordinary sentence case:
%
%   the table captions, IEEEtran.cls line 2778, which puts the caption text through \scshape;
%   the section headings, IEEEtran.cls line 5473, whose heading font ends in \scshape.
%
% Those are the two the responsible person asked for on 2026-09-18, and they are the only two that
% are live in this class option. The other \scshape sites in the class, lines 2737, 5468, 5501 and
% 5790, belong to the conference and computer society options, which these papers do not use, and
% the figure captions never had it: figure captions were already sentence case while table captions
% were not, so this change also makes the two agree, which they did not before.
%
% What is deliberately NOT changed: the drop cap at the start of each section, IEEEtran.cls lines
% 5835 and 5841, which sets the large initial and the rest of its first word through
% \MakeUppercase. A large initial followed by a word in capitals is a typographic convention and
% not shouting, and it is a bigger aesthetic move than was asked for. One line in this file takes
% it off if it is ever wanted: \renewcommand{\IEEEPARstartWORDCAPSTYLE}{\relax}.
%
% The class file is NOT edited. It is a third party class, vendored in three trees, and a modified
% copy under the same name is what makes a build unreproducible for somebody else. etoolbox patches
% the two macros from the preamble instead, so one change here reaches all three papers exactly as
% the draft switch does.
%
% Both patches fail the build if they do not apply. That matters more than it looks: if the class
% is ever updated and the macro no longer matches, a silent failure would bring the capitals back
% and nobody would know which build changed them.
\usepackage{etoolbox}
\makeatletter
\patchcmd{\@makecaption}{\scshape}{}{}{%
  \PackageError{sidestripe}{Cannot take the small capitals out of the table captions}
    {IEEEtran's \string\@makecaption no longer contains \string\scshape where it did at line
     2778. Read the class and fix the patch rather than removing it.}}
\patchcmd{\section}{\scshape}{}{}{%
  \PackageError{sidestripe}{Cannot take the small capitals out of the section headings}
    {IEEEtran's \string\section no longer contains \string\scshape where it did at line 5473.
     Read the class and fix the patch rather than removing it.}}
\makeatother

% Mechanics of the stripe. eso-pic's \AddToShipoutPictureBG* draws in the background layer of the
% current page only, which is what keeps the body untouched: the text block is not moved by a
% single point, because nothing is added to the galley. \rotatebox comes from graphicx, which the
% papers already load for their figures, and the grey is xcolor. tikz with `remember picture,
% overlay' would do the same thing at the price of a second compilation pass and a much larger
% dependency, and a raw \rotatebox in a zero size box inside the body would have to be placed by
% hand relative to the text and would move with it.
%
% Geometry. The stripe sits \SideStripeX from the left edge of the paper and starts \SideStripeY
% from the bottom, reading upwards. Those are far enough inside the paper that a printer's trim
% does not eat them and far enough from the text block that they cannot collide with it; both are
% named here so that the three papers cannot drift apart.
% If the paper that inputs this file has no published-date.tex, say so in words rather than
% letting TeX report an undefined control sequence from inside a rotated box. The three papers of
% paper/00001 to 00003 are in exactly that state on 2026-09-22 and are not rebuilt today; whoever
% builds one next is told here what to add and why, instead of debugging it.
\ifdefined\PublishedStripe\else
  \errmessage{sidestripe.tex: this paper inputs the stripe but defines no \string\PublishedStripe.
  Add a published-date.tex beside the paper and \string% The day this folder entered the PUBLIC repository, written here by hand once and never taken
% from the clock: a build that read the date from the machine would move it every time and a
% front page would claim a day it was not published on. The last updated date beside it in the
% stripe is a different quantity and comes from SOURCE_DATE_EPOCH.
%
% This folder is NOT in the public repository yet: on 2026-09-22 the only folder under
% results/ of evilaliv3/vesuvius-challenge-pipeline is e89eb7ad-umbilicus-estimator. So the
% stripe carries «not yet published» rather than a date, and when the folder is published the
% two lines below are replaced with the day it was, in this shape:
%
%   \newcommand{\PublishedDateIso}{2026-09-22}
%   \newcommand{\PublishedStripe}{published 22 September 2026}
%
% build.sh reads \PublishedDateIso and refuses a date in the future.
\newcommand{\PublishedDateIso}{}
\newcommand{\PublishedStripe}{not yet published}
. It holds the day
  the folder entered the PUBLIC repository, written by hand and never from the clock; a folder not
  yet published defines \string\PublishedStripe as: not yet published. See REVISION.md,
  2026-09-22}
\fi

\newcommand{\sidestripe}[1]{%
  \AddToShipoutPictureBG*{%
    \AtPageLowerLeft{%
      \put(\LenToUnit{\SideStripeX},\LenToUnit{\SideStripeY}){%
        \rotatebox{90}{%
          \color[gray]{0.55}%
          \fontsize{7}{8}\selectfont\ttfamily
          \ifSeriesDraft\SeriesDraftWord\SideStripeSep\fi%
          \SideStripeRepo\SideStripeSep results/#1\SideStripeSep
          \PublishedStripe\SideStripeSep last updated \BuildDateLong}}}}}

% CHANGE OF 2026-09-22, on the owner's word, to the one file all the papers of the series share.
%
% The stripe now carries TWO dates and names the folder as results/ rather than research/:
%
%   DRAFT . evilaliv3/vesuvius-challenge-pipeline . results/<folder> . <published> . last updated <date>
%
% \PublishedStripe comes from a published-date.tex beside the paper and is the day the folder
% entered the PUBLIC repository, written by hand once and never taken from the clock; a folder
% that is not published yet defines it as «not yet published». \BuildDateLong is the other
% quantity and keeps coming from SOURCE_DATE_EPOCH. A build refuses a publication date that is
% after the build date.
%
% EVERY PAPER THAT INPUTS THIS FILE MUST NOW ALSO INPUT A published-date.tex, or \PublishedStripe
% is undefined and the build stops. The three published papers of paper/00001 to 00003 are NOT
% rebuilt today: they are recompiled with their true publication dates at their next revision,
% which is the owner's decision of 2026-09-22 and is written in REVISION.md with its date.
, and the build of each one refuses a paper that does not, so none of them can
% miss the switch and no paper of the series can be marked while another is not.
%
% Off, this file draws no watermark and loads no package it did not load before, and each paper
% rebuilds to the bytes it had before the mark existed. That was checked on 2026-09-18 by building
% all three with the switch off and comparing the sha256 with the published files, not assumed.
\newif\ifSeriesDraft
\SeriesDraftfalse

\usepackage{eso-pic}
\usepackage{xcolor}

% The build writes build-date.tex and every paper of the series inputs it before this file. If it
% is missing, the date is not declared anywhere, and the build stops here rather than reaching a
% page that quietly carries no date or, worse, one taken from the clock.
\ifdefined\BuildDateLong\else
  \PackageError{sidestripe}{The command BuildDateLong is not defined: input build-date first}
    {The build writes build-date.tex from SOURCE_DATE_EPOCH. Run the paper's build script.}
\fi

\newlength{\SideStripeX}\setlength{\SideStripeX}{0.42in}
\newlength{\SideStripeY}\setlength{\SideStripeY}{1.05in}
% The separator is a middle dot: the house rule forbids a dash as punctuation, and a comma would
% read as a list rather than as the three fields of an identifier.
\newcommand{\SideStripeSep}{\,\textperiodcentered\,}
\newcommand{\SideStripeRepo}{evilaliv3/vesuvius-challenge-pipeline}

% The watermark, drawn only while the switch is on.
%
% Why a mark on every page and not a banner on the first. A reader almost never meets these papers
% whole: a page is printed, a screenshot is pasted into a message, a section is forwarded. A mark
% that lives on the first page alone is gone at that moment, and the page that travels is the one
% that will be quoted back. A running head would reach every page too, but it belongs to the
% galley: in this two column class it would move the text block, change where the columns break
% and cost a page, and it would have to be removed from the galley again to get the final paper
% back. A shipout picture costs nothing of the kind.
%
% Why the FOREGROUND and not the background. It was drawn in the background until 2026-09-18, and
% there it was invisible wherever a figure sat, because every figure in these papers fills its own
% box with opaque white before it draws anything. A figure is exactly the part of a page somebody
% screenshots, so the mark was missing from the part most likely to travel alone. On top of
% everything it reaches the whole page, figures included.
%
% What keeps it from spoiling the page. It is drawn at \SeriesDraftOpacity of full opacity, so what
% lands on the paper is a tint and not a shape: over white it is the lightest grey a PDF can carry,
% and over black text it changes nothing at all, because a tint of black over black is black. The
% text underneath is therefore not dimmed in any way a reader can measure. Nothing is added to the
% galley either, so the page count, the line breaks and the column breaks are exactly those of the
% final paper, and the gate on overfull boxes sees nothing new.
%
% The colour is black and not a grey, because at this opacity the colour is the only thing left to
% spend: the tint over white is 1 minus the opacity of the way from white to the ink, so black is
% the darkest mark two per cent can buy. A light grey at two per cent would be indistinguishable
% from the paper.
%
% \resizebox scales a box set at the normal size instead of asking for a font at 80pt, which the
% Computer Modern shapes of this class do not have: LaTeX would substitute the nearest design size
% and print a mark several times smaller than the one asked for, with a font warning the gates do
% not read. Scaling the box cannot fail that way, and it is a linear transform in the PDF, so the
% result is the same bytes on every build.
\ifSeriesDraft
  % transparent.sty writes the constant alpha into an ExtGState, which is what pdfTeX needs to put
  % ink on top of ink without hiding it. It is loaded only here, so a paper built with the switch
  % off carries no transparency machinery and no extra resource at all.
  \usepackage{transparent}
  \newcommand{\SeriesDraftWord}{DRAFT}
  \newcommand{\SeriesDraftOpacity}{0.10}
  \AddToShipoutPictureFG{%
    \AtPageCenter{%
      \makebox(0,0){%
        \rotatebox{45}{%
          \transparent{\SeriesDraftOpacity}%
          \resizebox{0.68\paperwidth}{!}{%
            \color{black}\sffamily\bfseries\SeriesDraftWord}}}}}
\fi

% THE CASE OF THE SERIES: the headings and the table captions read as the rest of the article.
%
% IEEEtran sets two things in small capitals, and a reader sees them as ALL CAPITALS shouted in the
% middle of a page whose every other line is ordinary sentence case:
%
%   the table captions, IEEEtran.cls line 2778, which puts the caption text through \scshape;
%   the section headings, IEEEtran.cls line 5473, whose heading font ends in \scshape.
%
% Those are the two the responsible person asked for on 2026-09-18, and they are the only two that
% are live in this class option. The other \scshape sites in the class, lines 2737, 5468, 5501 and
% 5790, belong to the conference and computer society options, which these papers do not use, and
% the figure captions never had it: figure captions were already sentence case while table captions
% were not, so this change also makes the two agree, which they did not before.
%
% What is deliberately NOT changed: the drop cap at the start of each section, IEEEtran.cls lines
% 5835 and 5841, which sets the large initial and the rest of its first word through
% \MakeUppercase. A large initial followed by a word in capitals is a typographic convention and
% not shouting, and it is a bigger aesthetic move than was asked for. One line in this file takes
% it off if it is ever wanted: \renewcommand{\IEEEPARstartWORDCAPSTYLE}{\relax}.
%
% The class file is NOT edited. It is a third party class, vendored in three trees, and a modified
% copy under the same name is what makes a build unreproducible for somebody else. etoolbox patches
% the two macros from the preamble instead, so one change here reaches all three papers exactly as
% the draft switch does.
%
% Both patches fail the build if they do not apply. That matters more than it looks: if the class
% is ever updated and the macro no longer matches, a silent failure would bring the capitals back
% and nobody would know which build changed them.
\usepackage{etoolbox}
\makeatletter
\patchcmd{\@makecaption}{\scshape}{}{}{%
  \PackageError{sidestripe}{Cannot take the small capitals out of the table captions}
    {IEEEtran's \string\@makecaption no longer contains \string\scshape where it did at line
     2778. Read the class and fix the patch rather than removing it.}}
\patchcmd{\section}{\scshape}{}{}{%
  \PackageError{sidestripe}{Cannot take the small capitals out of the section headings}
    {IEEEtran's \string\section no longer contains \string\scshape where it did at line 5473.
     Read the class and fix the patch rather than removing it.}}
\makeatother

% Mechanics of the stripe. eso-pic's \AddToShipoutPictureBG* draws in the background layer of the
% current page only, which is what keeps the body untouched: the text block is not moved by a
% single point, because nothing is added to the galley. \rotatebox comes from graphicx, which the
% papers already load for their figures, and the grey is xcolor. tikz with `remember picture,
% overlay' would do the same thing at the price of a second compilation pass and a much larger
% dependency, and a raw \rotatebox in a zero size box inside the body would have to be placed by
% hand relative to the text and would move with it.
%
% Geometry. The stripe sits \SideStripeX from the left edge of the paper and starts \SideStripeY
% from the bottom, reading upwards. Those are far enough inside the paper that a printer's trim
% does not eat them and far enough from the text block that they cannot collide with it; both are
% named here so that the three papers cannot drift apart.
% If the paper that inputs this file has no published-date.tex, say so in words rather than
% letting TeX report an undefined control sequence from inside a rotated box. The three papers of
% paper/00001 to 00003 are in exactly that state on 2026-09-22 and are not rebuilt today; whoever
% builds one next is told here what to add and why, instead of debugging it.
\ifdefined\PublishedStripe\else
  \errmessage{sidestripe.tex: this paper inputs the stripe but defines no \string\PublishedStripe.
  Add a published-date.tex beside the paper and \string% The day this folder entered the PUBLIC repository, written here by hand once and never taken
% from the clock: a build that read the date from the machine would move it every time and a
% front page would claim a day it was not published on. The last updated date beside it in the
% stripe is a different quantity and comes from SOURCE_DATE_EPOCH.
%
% This folder is NOT in the public repository yet: on 2026-09-22 the only folder under
% results/ of evilaliv3/vesuvius-challenge-pipeline is e89eb7ad-umbilicus-estimator. So the
% stripe carries «not yet published» rather than a date, and when the folder is published the
% two lines below are replaced with the day it was, in this shape:
%
%   \newcommand{\PublishedDateIso}{2026-09-22}
%   \newcommand{\PublishedStripe}{published 22 September 2026}
%
% build.sh reads \PublishedDateIso and refuses a date in the future.
\newcommand{\PublishedDateIso}{}
\newcommand{\PublishedStripe}{not yet published}
. It holds the day
  the folder entered the PUBLIC repository, written by hand and never from the clock; a folder not
  yet published defines \string\PublishedStripe as: not yet published. See REVISION.md,
  2026-09-22}
\fi

\newcommand{\sidestripe}[1]{%
  \AddToShipoutPictureBG*{%
    \AtPageLowerLeft{%
      \put(\LenToUnit{\SideStripeX},\LenToUnit{\SideStripeY}){%
        \rotatebox{90}{%
          \color[gray]{0.55}%
          \fontsize{7}{8}\selectfont\ttfamily
          \ifSeriesDraft\SeriesDraftWord\SideStripeSep\fi%
          \SideStripeRepo\SideStripeSep results/#1\SideStripeSep
          \PublishedStripe\SideStripeSep last updated \BuildDateLong}}}}}

% CHANGE OF 2026-09-22, on the owner's word, to the one file all the papers of the series share.
%
% The stripe now carries TWO dates and names the folder as results/ rather than research/:
%
%   DRAFT . evilaliv3/vesuvius-challenge-pipeline . results/<folder> . <published> . last updated <date>
%
% \PublishedStripe comes from a published-date.tex beside the paper and is the day the folder
% entered the PUBLIC repository, written by hand once and never taken from the clock; a folder
% that is not published yet defines it as «not yet published». \BuildDateLong is the other
% quantity and keeps coming from SOURCE_DATE_EPOCH. A build refuses a publication date that is
% after the build date.
%
% EVERY PAPER THAT INPUTS THIS FILE MUST NOW ALSO INPUT A published-date.tex, or \PublishedStripe
% is undefined and the build stops. The three published papers of paper/00001 to 00003 are NOT
% rebuilt today: they are recompiled with their true publication dates at their next revision,
% which is the owner's decision of 2026-09-22 and is written in REVISION.md with its date.

%
% and in the BODY, immediately after \begin{document}, with the paper's own folder number:
%
%     \sidestripe{00001}
%
% The second line has to be in the body because the stripe is attached to the next page shipped
% out, which is the first page; the first line has to be in the preamble because it loads packages.
% The build of each paper fails if it carries neither.
%
% The date is \BuildDateLong, written by the build from SOURCE_DATE_EPOCH, the same value pdfTeX
% stamps into /CreationDate. It is never typed, so the front page and the metadata cannot come to
% say two different days. It is the date of the revision and not a reading of the clock, which is
% why the stripe labels it "last updated": a bare date beside a title is read as the day of
% writing, of submission or of the build, and a reader has no way to tell which was meant.

% THE DRAFT SWITCH OF THE SERIES. This one line is the whole control, for all three papers.
%
%     \SeriesDrafttrue    the papers are drafts and carry the mark
%     \SeriesDraftfalse   the papers are final and carry nothing
%
% Nothing else moves when it is thrown. The three papers reach this file through
% % The marks of the series: what a paper is, in the left margin of its first page, and, while the
% series is a draft, the word DRAFT across every page of it.
%
% One file for the three papers, so that a later change to the wording, the colour or the position
% lands in all three at once. It draws, rotated a quarter turn and read from the bottom, the state
% of the paper, the repository, the paper's own folder and the date of the revision it prints:
%
%     DRAFT . evilaliv3/vesuvius-challenge-pipeline . research/00001 . last updated 18 September 2026
%
% Why it exists. These are self published reports whose claims are pinned to a codebase that moves
% under them and to data that changes, and the method they argue for rests on when a thing was
% declared against when it was measured. A paper that hides its own date is inconsistent with the
% way it argues, and a reader who prints the front page should be able to see, without reading a
% line of it, what the thing is, what state it is in and when it was last revised.
%
% Usage, two lines per paper. In the PREAMBLE, after the class and before \begin{document}:
%
%     % The marks of the series: what a paper is, in the left margin of its first page, and, while the
% series is a draft, the word DRAFT across every page of it.
%
% One file for the three papers, so that a later change to the wording, the colour or the position
% lands in all three at once. It draws, rotated a quarter turn and read from the bottom, the state
% of the paper, the repository, the paper's own folder and the date of the revision it prints:
%
%     DRAFT . evilaliv3/vesuvius-challenge-pipeline . research/00001 . last updated 18 September 2026
%
% Why it exists. These are self published reports whose claims are pinned to a codebase that moves
% under them and to data that changes, and the method they argue for rests on when a thing was
% declared against when it was measured. A paper that hides its own date is inconsistent with the
% way it argues, and a reader who prints the front page should be able to see, without reading a
% line of it, what the thing is, what state it is in and when it was last revised.
%
% Usage, two lines per paper. In the PREAMBLE, after the class and before \begin{document}:
%
%     % The marks of the series: what a paper is, in the left margin of its first page, and, while the
% series is a draft, the word DRAFT across every page of it.
%
% One file for the three papers, so that a later change to the wording, the colour or the position
% lands in all three at once. It draws, rotated a quarter turn and read from the bottom, the state
% of the paper, the repository, the paper's own folder and the date of the revision it prints:
%
%     DRAFT . evilaliv3/vesuvius-challenge-pipeline . research/00001 . last updated 18 September 2026
%
% Why it exists. These are self published reports whose claims are pinned to a codebase that moves
% under them and to data that changes, and the method they argue for rests on when a thing was
% declared against when it was measured. A paper that hides its own date is inconsistent with the
% way it argues, and a reader who prints the front page should be able to see, without reading a
% line of it, what the thing is, what state it is in and when it was last revised.
%
% Usage, two lines per paper. In the PREAMBLE, after the class and before \begin{document}:
%
%     \input{sidestripe}
%
% and in the BODY, immediately after \begin{document}, with the paper's own folder number:
%
%     \sidestripe{00001}
%
% The second line has to be in the body because the stripe is attached to the next page shipped
% out, which is the first page; the first line has to be in the preamble because it loads packages.
% The build of each paper fails if it carries neither.
%
% The date is \BuildDateLong, written by the build from SOURCE_DATE_EPOCH, the same value pdfTeX
% stamps into /CreationDate. It is never typed, so the front page and the metadata cannot come to
% say two different days. It is the date of the revision and not a reading of the clock, which is
% why the stripe labels it "last updated": a bare date beside a title is read as the day of
% writing, of submission or of the build, and a reader has no way to tell which was meant.

% THE DRAFT SWITCH OF THE SERIES. This one line is the whole control, for all three papers.
%
%     \SeriesDrafttrue    the papers are drafts and carry the mark
%     \SeriesDraftfalse   the papers are final and carry nothing
%
% Nothing else moves when it is thrown. The three papers reach this file through
% \input{sidestripe}, and the build of each one refuses a paper that does not, so none of them can
% miss the switch and no paper of the series can be marked while another is not.
%
% Off, this file draws no watermark and loads no package it did not load before, and each paper
% rebuilds to the bytes it had before the mark existed. That was checked on 2026-09-18 by building
% all three with the switch off and comparing the sha256 with the published files, not assumed.
\newif\ifSeriesDraft
\SeriesDraftfalse

\usepackage{eso-pic}
\usepackage{xcolor}

% The build writes build-date.tex and every paper of the series inputs it before this file. If it
% is missing, the date is not declared anywhere, and the build stops here rather than reaching a
% page that quietly carries no date or, worse, one taken from the clock.
\ifdefined\BuildDateLong\else
  \PackageError{sidestripe}{The command BuildDateLong is not defined: input build-date first}
    {The build writes build-date.tex from SOURCE_DATE_EPOCH. Run the paper's build script.}
\fi

\newlength{\SideStripeX}\setlength{\SideStripeX}{0.42in}
\newlength{\SideStripeY}\setlength{\SideStripeY}{1.05in}
% The separator is a middle dot: the house rule forbids a dash as punctuation, and a comma would
% read as a list rather than as the three fields of an identifier.
\newcommand{\SideStripeSep}{\,\textperiodcentered\,}
\newcommand{\SideStripeRepo}{evilaliv3/vesuvius-challenge-pipeline}

% The watermark, drawn only while the switch is on.
%
% Why a mark on every page and not a banner on the first. A reader almost never meets these papers
% whole: a page is printed, a screenshot is pasted into a message, a section is forwarded. A mark
% that lives on the first page alone is gone at that moment, and the page that travels is the one
% that will be quoted back. A running head would reach every page too, but it belongs to the
% galley: in this two column class it would move the text block, change where the columns break
% and cost a page, and it would have to be removed from the galley again to get the final paper
% back. A shipout picture costs nothing of the kind.
%
% Why the FOREGROUND and not the background. It was drawn in the background until 2026-09-18, and
% there it was invisible wherever a figure sat, because every figure in these papers fills its own
% box with opaque white before it draws anything. A figure is exactly the part of a page somebody
% screenshots, so the mark was missing from the part most likely to travel alone. On top of
% everything it reaches the whole page, figures included.
%
% What keeps it from spoiling the page. It is drawn at \SeriesDraftOpacity of full opacity, so what
% lands on the paper is a tint and not a shape: over white it is the lightest grey a PDF can carry,
% and over black text it changes nothing at all, because a tint of black over black is black. The
% text underneath is therefore not dimmed in any way a reader can measure. Nothing is added to the
% galley either, so the page count, the line breaks and the column breaks are exactly those of the
% final paper, and the gate on overfull boxes sees nothing new.
%
% The colour is black and not a grey, because at this opacity the colour is the only thing left to
% spend: the tint over white is 1 minus the opacity of the way from white to the ink, so black is
% the darkest mark two per cent can buy. A light grey at two per cent would be indistinguishable
% from the paper.
%
% \resizebox scales a box set at the normal size instead of asking for a font at 80pt, which the
% Computer Modern shapes of this class do not have: LaTeX would substitute the nearest design size
% and print a mark several times smaller than the one asked for, with a font warning the gates do
% not read. Scaling the box cannot fail that way, and it is a linear transform in the PDF, so the
% result is the same bytes on every build.
\ifSeriesDraft
  % transparent.sty writes the constant alpha into an ExtGState, which is what pdfTeX needs to put
  % ink on top of ink without hiding it. It is loaded only here, so a paper built with the switch
  % off carries no transparency machinery and no extra resource at all.
  \usepackage{transparent}
  \newcommand{\SeriesDraftWord}{DRAFT}
  \newcommand{\SeriesDraftOpacity}{0.10}
  \AddToShipoutPictureFG{%
    \AtPageCenter{%
      \makebox(0,0){%
        \rotatebox{45}{%
          \transparent{\SeriesDraftOpacity}%
          \resizebox{0.68\paperwidth}{!}{%
            \color{black}\sffamily\bfseries\SeriesDraftWord}}}}}
\fi

% THE CASE OF THE SERIES: the headings and the table captions read as the rest of the article.
%
% IEEEtran sets two things in small capitals, and a reader sees them as ALL CAPITALS shouted in the
% middle of a page whose every other line is ordinary sentence case:
%
%   the table captions, IEEEtran.cls line 2778, which puts the caption text through \scshape;
%   the section headings, IEEEtran.cls line 5473, whose heading font ends in \scshape.
%
% Those are the two the responsible person asked for on 2026-09-18, and they are the only two that
% are live in this class option. The other \scshape sites in the class, lines 2737, 5468, 5501 and
% 5790, belong to the conference and computer society options, which these papers do not use, and
% the figure captions never had it: figure captions were already sentence case while table captions
% were not, so this change also makes the two agree, which they did not before.
%
% What is deliberately NOT changed: the drop cap at the start of each section, IEEEtran.cls lines
% 5835 and 5841, which sets the large initial and the rest of its first word through
% \MakeUppercase. A large initial followed by a word in capitals is a typographic convention and
% not shouting, and it is a bigger aesthetic move than was asked for. One line in this file takes
% it off if it is ever wanted: \renewcommand{\IEEEPARstartWORDCAPSTYLE}{\relax}.
%
% The class file is NOT edited. It is a third party class, vendored in three trees, and a modified
% copy under the same name is what makes a build unreproducible for somebody else. etoolbox patches
% the two macros from the preamble instead, so one change here reaches all three papers exactly as
% the draft switch does.
%
% Both patches fail the build if they do not apply. That matters more than it looks: if the class
% is ever updated and the macro no longer matches, a silent failure would bring the capitals back
% and nobody would know which build changed them.
\usepackage{etoolbox}
\makeatletter
\patchcmd{\@makecaption}{\scshape}{}{}{%
  \PackageError{sidestripe}{Cannot take the small capitals out of the table captions}
    {IEEEtran's \string\@makecaption no longer contains \string\scshape where it did at line
     2778. Read the class and fix the patch rather than removing it.}}
\patchcmd{\section}{\scshape}{}{}{%
  \PackageError{sidestripe}{Cannot take the small capitals out of the section headings}
    {IEEEtran's \string\section no longer contains \string\scshape where it did at line 5473.
     Read the class and fix the patch rather than removing it.}}
\makeatother

% Mechanics of the stripe. eso-pic's \AddToShipoutPictureBG* draws in the background layer of the
% current page only, which is what keeps the body untouched: the text block is not moved by a
% single point, because nothing is added to the galley. \rotatebox comes from graphicx, which the
% papers already load for their figures, and the grey is xcolor. tikz with `remember picture,
% overlay' would do the same thing at the price of a second compilation pass and a much larger
% dependency, and a raw \rotatebox in a zero size box inside the body would have to be placed by
% hand relative to the text and would move with it.
%
% Geometry. The stripe sits \SideStripeX from the left edge of the paper and starts \SideStripeY
% from the bottom, reading upwards. Those are far enough inside the paper that a printer's trim
% does not eat them and far enough from the text block that they cannot collide with it; both are
% named here so that the three papers cannot drift apart.
% If the paper that inputs this file has no published-date.tex, say so in words rather than
% letting TeX report an undefined control sequence from inside a rotated box. The three papers of
% paper/00001 to 00003 are in exactly that state on 2026-09-22 and are not rebuilt today; whoever
% builds one next is told here what to add and why, instead of debugging it.
\ifdefined\PublishedStripe\else
  \errmessage{sidestripe.tex: this paper inputs the stripe but defines no \string\PublishedStripe.
  Add a published-date.tex beside the paper and \string\input{published-date}. It holds the day
  the folder entered the PUBLIC repository, written by hand and never from the clock; a folder not
  yet published defines \string\PublishedStripe as: not yet published. See REVISION.md,
  2026-09-22}
\fi

\newcommand{\sidestripe}[1]{%
  \AddToShipoutPictureBG*{%
    \AtPageLowerLeft{%
      \put(\LenToUnit{\SideStripeX},\LenToUnit{\SideStripeY}){%
        \rotatebox{90}{%
          \color[gray]{0.55}%
          \fontsize{7}{8}\selectfont\ttfamily
          \ifSeriesDraft\SeriesDraftWord\SideStripeSep\fi%
          \SideStripeRepo\SideStripeSep results/#1\SideStripeSep
          \PublishedStripe\SideStripeSep last updated \BuildDateLong}}}}}

% CHANGE OF 2026-09-22, on the owner's word, to the one file all the papers of the series share.
%
% The stripe now carries TWO dates and names the folder as results/ rather than research/:
%
%   DRAFT . evilaliv3/vesuvius-challenge-pipeline . results/<folder> . <published> . last updated <date>
%
% \PublishedStripe comes from a published-date.tex beside the paper and is the day the folder
% entered the PUBLIC repository, written by hand once and never taken from the clock; a folder
% that is not published yet defines it as «not yet published». \BuildDateLong is the other
% quantity and keeps coming from SOURCE_DATE_EPOCH. A build refuses a publication date that is
% after the build date.
%
% EVERY PAPER THAT INPUTS THIS FILE MUST NOW ALSO INPUT A published-date.tex, or \PublishedStripe
% is undefined and the build stops. The three published papers of paper/00001 to 00003 are NOT
% rebuilt today: they are recompiled with their true publication dates at their next revision,
% which is the owner's decision of 2026-09-22 and is written in REVISION.md with its date.

%
% and in the BODY, immediately after \begin{document}, with the paper's own folder number:
%
%     \sidestripe{00001}
%
% The second line has to be in the body because the stripe is attached to the next page shipped
% out, which is the first page; the first line has to be in the preamble because it loads packages.
% The build of each paper fails if it carries neither.
%
% The date is \BuildDateLong, written by the build from SOURCE_DATE_EPOCH, the same value pdfTeX
% stamps into /CreationDate. It is never typed, so the front page and the metadata cannot come to
% say two different days. It is the date of the revision and not a reading of the clock, which is
% why the stripe labels it "last updated": a bare date beside a title is read as the day of
% writing, of submission or of the build, and a reader has no way to tell which was meant.

% THE DRAFT SWITCH OF THE SERIES. This one line is the whole control, for all three papers.
%
%     \SeriesDrafttrue    the papers are drafts and carry the mark
%     \SeriesDraftfalse   the papers are final and carry nothing
%
% Nothing else moves when it is thrown. The three papers reach this file through
% % The marks of the series: what a paper is, in the left margin of its first page, and, while the
% series is a draft, the word DRAFT across every page of it.
%
% One file for the three papers, so that a later change to the wording, the colour or the position
% lands in all three at once. It draws, rotated a quarter turn and read from the bottom, the state
% of the paper, the repository, the paper's own folder and the date of the revision it prints:
%
%     DRAFT . evilaliv3/vesuvius-challenge-pipeline . research/00001 . last updated 18 September 2026
%
% Why it exists. These are self published reports whose claims are pinned to a codebase that moves
% under them and to data that changes, and the method they argue for rests on when a thing was
% declared against when it was measured. A paper that hides its own date is inconsistent with the
% way it argues, and a reader who prints the front page should be able to see, without reading a
% line of it, what the thing is, what state it is in and when it was last revised.
%
% Usage, two lines per paper. In the PREAMBLE, after the class and before \begin{document}:
%
%     \input{sidestripe}
%
% and in the BODY, immediately after \begin{document}, with the paper's own folder number:
%
%     \sidestripe{00001}
%
% The second line has to be in the body because the stripe is attached to the next page shipped
% out, which is the first page; the first line has to be in the preamble because it loads packages.
% The build of each paper fails if it carries neither.
%
% The date is \BuildDateLong, written by the build from SOURCE_DATE_EPOCH, the same value pdfTeX
% stamps into /CreationDate. It is never typed, so the front page and the metadata cannot come to
% say two different days. It is the date of the revision and not a reading of the clock, which is
% why the stripe labels it "last updated": a bare date beside a title is read as the day of
% writing, of submission or of the build, and a reader has no way to tell which was meant.

% THE DRAFT SWITCH OF THE SERIES. This one line is the whole control, for all three papers.
%
%     \SeriesDrafttrue    the papers are drafts and carry the mark
%     \SeriesDraftfalse   the papers are final and carry nothing
%
% Nothing else moves when it is thrown. The three papers reach this file through
% \input{sidestripe}, and the build of each one refuses a paper that does not, so none of them can
% miss the switch and no paper of the series can be marked while another is not.
%
% Off, this file draws no watermark and loads no package it did not load before, and each paper
% rebuilds to the bytes it had before the mark existed. That was checked on 2026-09-18 by building
% all three with the switch off and comparing the sha256 with the published files, not assumed.
\newif\ifSeriesDraft
\SeriesDraftfalse

\usepackage{eso-pic}
\usepackage{xcolor}

% The build writes build-date.tex and every paper of the series inputs it before this file. If it
% is missing, the date is not declared anywhere, and the build stops here rather than reaching a
% page that quietly carries no date or, worse, one taken from the clock.
\ifdefined\BuildDateLong\else
  \PackageError{sidestripe}{The command BuildDateLong is not defined: input build-date first}
    {The build writes build-date.tex from SOURCE_DATE_EPOCH. Run the paper's build script.}
\fi

\newlength{\SideStripeX}\setlength{\SideStripeX}{0.42in}
\newlength{\SideStripeY}\setlength{\SideStripeY}{1.05in}
% The separator is a middle dot: the house rule forbids a dash as punctuation, and a comma would
% read as a list rather than as the three fields of an identifier.
\newcommand{\SideStripeSep}{\,\textperiodcentered\,}
\newcommand{\SideStripeRepo}{evilaliv3/vesuvius-challenge-pipeline}

% The watermark, drawn only while the switch is on.
%
% Why a mark on every page and not a banner on the first. A reader almost never meets these papers
% whole: a page is printed, a screenshot is pasted into a message, a section is forwarded. A mark
% that lives on the first page alone is gone at that moment, and the page that travels is the one
% that will be quoted back. A running head would reach every page too, but it belongs to the
% galley: in this two column class it would move the text block, change where the columns break
% and cost a page, and it would have to be removed from the galley again to get the final paper
% back. A shipout picture costs nothing of the kind.
%
% Why the FOREGROUND and not the background. It was drawn in the background until 2026-09-18, and
% there it was invisible wherever a figure sat, because every figure in these papers fills its own
% box with opaque white before it draws anything. A figure is exactly the part of a page somebody
% screenshots, so the mark was missing from the part most likely to travel alone. On top of
% everything it reaches the whole page, figures included.
%
% What keeps it from spoiling the page. It is drawn at \SeriesDraftOpacity of full opacity, so what
% lands on the paper is a tint and not a shape: over white it is the lightest grey a PDF can carry,
% and over black text it changes nothing at all, because a tint of black over black is black. The
% text underneath is therefore not dimmed in any way a reader can measure. Nothing is added to the
% galley either, so the page count, the line breaks and the column breaks are exactly those of the
% final paper, and the gate on overfull boxes sees nothing new.
%
% The colour is black and not a grey, because at this opacity the colour is the only thing left to
% spend: the tint over white is 1 minus the opacity of the way from white to the ink, so black is
% the darkest mark two per cent can buy. A light grey at two per cent would be indistinguishable
% from the paper.
%
% \resizebox scales a box set at the normal size instead of asking for a font at 80pt, which the
% Computer Modern shapes of this class do not have: LaTeX would substitute the nearest design size
% and print a mark several times smaller than the one asked for, with a font warning the gates do
% not read. Scaling the box cannot fail that way, and it is a linear transform in the PDF, so the
% result is the same bytes on every build.
\ifSeriesDraft
  % transparent.sty writes the constant alpha into an ExtGState, which is what pdfTeX needs to put
  % ink on top of ink without hiding it. It is loaded only here, so a paper built with the switch
  % off carries no transparency machinery and no extra resource at all.
  \usepackage{transparent}
  \newcommand{\SeriesDraftWord}{DRAFT}
  \newcommand{\SeriesDraftOpacity}{0.10}
  \AddToShipoutPictureFG{%
    \AtPageCenter{%
      \makebox(0,0){%
        \rotatebox{45}{%
          \transparent{\SeriesDraftOpacity}%
          \resizebox{0.68\paperwidth}{!}{%
            \color{black}\sffamily\bfseries\SeriesDraftWord}}}}}
\fi

% THE CASE OF THE SERIES: the headings and the table captions read as the rest of the article.
%
% IEEEtran sets two things in small capitals, and a reader sees them as ALL CAPITALS shouted in the
% middle of a page whose every other line is ordinary sentence case:
%
%   the table captions, IEEEtran.cls line 2778, which puts the caption text through \scshape;
%   the section headings, IEEEtran.cls line 5473, whose heading font ends in \scshape.
%
% Those are the two the responsible person asked for on 2026-09-18, and they are the only two that
% are live in this class option. The other \scshape sites in the class, lines 2737, 5468, 5501 and
% 5790, belong to the conference and computer society options, which these papers do not use, and
% the figure captions never had it: figure captions were already sentence case while table captions
% were not, so this change also makes the two agree, which they did not before.
%
% What is deliberately NOT changed: the drop cap at the start of each section, IEEEtran.cls lines
% 5835 and 5841, which sets the large initial and the rest of its first word through
% \MakeUppercase. A large initial followed by a word in capitals is a typographic convention and
% not shouting, and it is a bigger aesthetic move than was asked for. One line in this file takes
% it off if it is ever wanted: \renewcommand{\IEEEPARstartWORDCAPSTYLE}{\relax}.
%
% The class file is NOT edited. It is a third party class, vendored in three trees, and a modified
% copy under the same name is what makes a build unreproducible for somebody else. etoolbox patches
% the two macros from the preamble instead, so one change here reaches all three papers exactly as
% the draft switch does.
%
% Both patches fail the build if they do not apply. That matters more than it looks: if the class
% is ever updated and the macro no longer matches, a silent failure would bring the capitals back
% and nobody would know which build changed them.
\usepackage{etoolbox}
\makeatletter
\patchcmd{\@makecaption}{\scshape}{}{}{%
  \PackageError{sidestripe}{Cannot take the small capitals out of the table captions}
    {IEEEtran's \string\@makecaption no longer contains \string\scshape where it did at line
     2778. Read the class and fix the patch rather than removing it.}}
\patchcmd{\section}{\scshape}{}{}{%
  \PackageError{sidestripe}{Cannot take the small capitals out of the section headings}
    {IEEEtran's \string\section no longer contains \string\scshape where it did at line 5473.
     Read the class and fix the patch rather than removing it.}}
\makeatother

% Mechanics of the stripe. eso-pic's \AddToShipoutPictureBG* draws in the background layer of the
% current page only, which is what keeps the body untouched: the text block is not moved by a
% single point, because nothing is added to the galley. \rotatebox comes from graphicx, which the
% papers already load for their figures, and the grey is xcolor. tikz with `remember picture,
% overlay' would do the same thing at the price of a second compilation pass and a much larger
% dependency, and a raw \rotatebox in a zero size box inside the body would have to be placed by
% hand relative to the text and would move with it.
%
% Geometry. The stripe sits \SideStripeX from the left edge of the paper and starts \SideStripeY
% from the bottom, reading upwards. Those are far enough inside the paper that a printer's trim
% does not eat them and far enough from the text block that they cannot collide with it; both are
% named here so that the three papers cannot drift apart.
% If the paper that inputs this file has no published-date.tex, say so in words rather than
% letting TeX report an undefined control sequence from inside a rotated box. The three papers of
% paper/00001 to 00003 are in exactly that state on 2026-09-22 and are not rebuilt today; whoever
% builds one next is told here what to add and why, instead of debugging it.
\ifdefined\PublishedStripe\else
  \errmessage{sidestripe.tex: this paper inputs the stripe but defines no \string\PublishedStripe.
  Add a published-date.tex beside the paper and \string\input{published-date}. It holds the day
  the folder entered the PUBLIC repository, written by hand and never from the clock; a folder not
  yet published defines \string\PublishedStripe as: not yet published. See REVISION.md,
  2026-09-22}
\fi

\newcommand{\sidestripe}[1]{%
  \AddToShipoutPictureBG*{%
    \AtPageLowerLeft{%
      \put(\LenToUnit{\SideStripeX},\LenToUnit{\SideStripeY}){%
        \rotatebox{90}{%
          \color[gray]{0.55}%
          \fontsize{7}{8}\selectfont\ttfamily
          \ifSeriesDraft\SeriesDraftWord\SideStripeSep\fi%
          \SideStripeRepo\SideStripeSep results/#1\SideStripeSep
          \PublishedStripe\SideStripeSep last updated \BuildDateLong}}}}}

% CHANGE OF 2026-09-22, on the owner's word, to the one file all the papers of the series share.
%
% The stripe now carries TWO dates and names the folder as results/ rather than research/:
%
%   DRAFT . evilaliv3/vesuvius-challenge-pipeline . results/<folder> . <published> . last updated <date>
%
% \PublishedStripe comes from a published-date.tex beside the paper and is the day the folder
% entered the PUBLIC repository, written by hand once and never taken from the clock; a folder
% that is not published yet defines it as «not yet published». \BuildDateLong is the other
% quantity and keeps coming from SOURCE_DATE_EPOCH. A build refuses a publication date that is
% after the build date.
%
% EVERY PAPER THAT INPUTS THIS FILE MUST NOW ALSO INPUT A published-date.tex, or \PublishedStripe
% is undefined and the build stops. The three published papers of paper/00001 to 00003 are NOT
% rebuilt today: they are recompiled with their true publication dates at their next revision,
% which is the owner's decision of 2026-09-22 and is written in REVISION.md with its date.
, and the build of each one refuses a paper that does not, so none of them can
% miss the switch and no paper of the series can be marked while another is not.
%
% Off, this file draws no watermark and loads no package it did not load before, and each paper
% rebuilds to the bytes it had before the mark existed. That was checked on 2026-09-18 by building
% all three with the switch off and comparing the sha256 with the published files, not assumed.
\newif\ifSeriesDraft
\SeriesDraftfalse

\usepackage{eso-pic}
\usepackage{xcolor}

% The build writes build-date.tex and every paper of the series inputs it before this file. If it
% is missing, the date is not declared anywhere, and the build stops here rather than reaching a
% page that quietly carries no date or, worse, one taken from the clock.
\ifdefined\BuildDateLong\else
  \PackageError{sidestripe}{The command BuildDateLong is not defined: input build-date first}
    {The build writes build-date.tex from SOURCE_DATE_EPOCH. Run the paper's build script.}
\fi

\newlength{\SideStripeX}\setlength{\SideStripeX}{0.42in}
\newlength{\SideStripeY}\setlength{\SideStripeY}{1.05in}
% The separator is a middle dot: the house rule forbids a dash as punctuation, and a comma would
% read as a list rather than as the three fields of an identifier.
\newcommand{\SideStripeSep}{\,\textperiodcentered\,}
\newcommand{\SideStripeRepo}{evilaliv3/vesuvius-challenge-pipeline}

% The watermark, drawn only while the switch is on.
%
% Why a mark on every page and not a banner on the first. A reader almost never meets these papers
% whole: a page is printed, a screenshot is pasted into a message, a section is forwarded. A mark
% that lives on the first page alone is gone at that moment, and the page that travels is the one
% that will be quoted back. A running head would reach every page too, but it belongs to the
% galley: in this two column class it would move the text block, change where the columns break
% and cost a page, and it would have to be removed from the galley again to get the final paper
% back. A shipout picture costs nothing of the kind.
%
% Why the FOREGROUND and not the background. It was drawn in the background until 2026-09-18, and
% there it was invisible wherever a figure sat, because every figure in these papers fills its own
% box with opaque white before it draws anything. A figure is exactly the part of a page somebody
% screenshots, so the mark was missing from the part most likely to travel alone. On top of
% everything it reaches the whole page, figures included.
%
% What keeps it from spoiling the page. It is drawn at \SeriesDraftOpacity of full opacity, so what
% lands on the paper is a tint and not a shape: over white it is the lightest grey a PDF can carry,
% and over black text it changes nothing at all, because a tint of black over black is black. The
% text underneath is therefore not dimmed in any way a reader can measure. Nothing is added to the
% galley either, so the page count, the line breaks and the column breaks are exactly those of the
% final paper, and the gate on overfull boxes sees nothing new.
%
% The colour is black and not a grey, because at this opacity the colour is the only thing left to
% spend: the tint over white is 1 minus the opacity of the way from white to the ink, so black is
% the darkest mark two per cent can buy. A light grey at two per cent would be indistinguishable
% from the paper.
%
% \resizebox scales a box set at the normal size instead of asking for a font at 80pt, which the
% Computer Modern shapes of this class do not have: LaTeX would substitute the nearest design size
% and print a mark several times smaller than the one asked for, with a font warning the gates do
% not read. Scaling the box cannot fail that way, and it is a linear transform in the PDF, so the
% result is the same bytes on every build.
\ifSeriesDraft
  % transparent.sty writes the constant alpha into an ExtGState, which is what pdfTeX needs to put
  % ink on top of ink without hiding it. It is loaded only here, so a paper built with the switch
  % off carries no transparency machinery and no extra resource at all.
  \usepackage{transparent}
  \newcommand{\SeriesDraftWord}{DRAFT}
  \newcommand{\SeriesDraftOpacity}{0.10}
  \AddToShipoutPictureFG{%
    \AtPageCenter{%
      \makebox(0,0){%
        \rotatebox{45}{%
          \transparent{\SeriesDraftOpacity}%
          \resizebox{0.68\paperwidth}{!}{%
            \color{black}\sffamily\bfseries\SeriesDraftWord}}}}}
\fi

% THE CASE OF THE SERIES: the headings and the table captions read as the rest of the article.
%
% IEEEtran sets two things in small capitals, and a reader sees them as ALL CAPITALS shouted in the
% middle of a page whose every other line is ordinary sentence case:
%
%   the table captions, IEEEtran.cls line 2778, which puts the caption text through \scshape;
%   the section headings, IEEEtran.cls line 5473, whose heading font ends in \scshape.
%
% Those are the two the responsible person asked for on 2026-09-18, and they are the only two that
% are live in this class option. The other \scshape sites in the class, lines 2737, 5468, 5501 and
% 5790, belong to the conference and computer society options, which these papers do not use, and
% the figure captions never had it: figure captions were already sentence case while table captions
% were not, so this change also makes the two agree, which they did not before.
%
% What is deliberately NOT changed: the drop cap at the start of each section, IEEEtran.cls lines
% 5835 and 5841, which sets the large initial and the rest of its first word through
% \MakeUppercase. A large initial followed by a word in capitals is a typographic convention and
% not shouting, and it is a bigger aesthetic move than was asked for. One line in this file takes
% it off if it is ever wanted: \renewcommand{\IEEEPARstartWORDCAPSTYLE}{\relax}.
%
% The class file is NOT edited. It is a third party class, vendored in three trees, and a modified
% copy under the same name is what makes a build unreproducible for somebody else. etoolbox patches
% the two macros from the preamble instead, so one change here reaches all three papers exactly as
% the draft switch does.
%
% Both patches fail the build if they do not apply. That matters more than it looks: if the class
% is ever updated and the macro no longer matches, a silent failure would bring the capitals back
% and nobody would know which build changed them.
\usepackage{etoolbox}
\makeatletter
\patchcmd{\@makecaption}{\scshape}{}{}{%
  \PackageError{sidestripe}{Cannot take the small capitals out of the table captions}
    {IEEEtran's \string\@makecaption no longer contains \string\scshape where it did at line
     2778. Read the class and fix the patch rather than removing it.}}
\patchcmd{\section}{\scshape}{}{}{%
  \PackageError{sidestripe}{Cannot take the small capitals out of the section headings}
    {IEEEtran's \string\section no longer contains \string\scshape where it did at line 5473.
     Read the class and fix the patch rather than removing it.}}
\makeatother

% Mechanics of the stripe. eso-pic's \AddToShipoutPictureBG* draws in the background layer of the
% current page only, which is what keeps the body untouched: the text block is not moved by a
% single point, because nothing is added to the galley. \rotatebox comes from graphicx, which the
% papers already load for their figures, and the grey is xcolor. tikz with `remember picture,
% overlay' would do the same thing at the price of a second compilation pass and a much larger
% dependency, and a raw \rotatebox in a zero size box inside the body would have to be placed by
% hand relative to the text and would move with it.
%
% Geometry. The stripe sits \SideStripeX from the left edge of the paper and starts \SideStripeY
% from the bottom, reading upwards. Those are far enough inside the paper that a printer's trim
% does not eat them and far enough from the text block that they cannot collide with it; both are
% named here so that the three papers cannot drift apart.
% If the paper that inputs this file has no published-date.tex, say so in words rather than
% letting TeX report an undefined control sequence from inside a rotated box. The three papers of
% paper/00001 to 00003 are in exactly that state on 2026-09-22 and are not rebuilt today; whoever
% builds one next is told here what to add and why, instead of debugging it.
\ifdefined\PublishedStripe\else
  \errmessage{sidestripe.tex: this paper inputs the stripe but defines no \string\PublishedStripe.
  Add a published-date.tex beside the paper and \string% The day this folder entered the PUBLIC repository, written here by hand once and never taken
% from the clock: a build that read the date from the machine would move it every time and a
% front page would claim a day it was not published on. The last updated date beside it in the
% stripe is a different quantity and comes from SOURCE_DATE_EPOCH.
%
% This folder is NOT in the public repository yet: on 2026-09-22 the only folder under
% results/ of evilaliv3/vesuvius-challenge-pipeline is e89eb7ad-umbilicus-estimator. So the
% stripe carries «not yet published» rather than a date, and when the folder is published the
% two lines below are replaced with the day it was, in this shape:
%
%   \newcommand{\PublishedDateIso}{2026-09-22}
%   \newcommand{\PublishedStripe}{published 22 September 2026}
%
% build.sh reads \PublishedDateIso and refuses a date in the future.
\newcommand{\PublishedDateIso}{}
\newcommand{\PublishedStripe}{not yet published}
. It holds the day
  the folder entered the PUBLIC repository, written by hand and never from the clock; a folder not
  yet published defines \string\PublishedStripe as: not yet published. See REVISION.md,
  2026-09-22}
\fi

\newcommand{\sidestripe}[1]{%
  \AddToShipoutPictureBG*{%
    \AtPageLowerLeft{%
      \put(\LenToUnit{\SideStripeX},\LenToUnit{\SideStripeY}){%
        \rotatebox{90}{%
          \color[gray]{0.55}%
          \fontsize{7}{8}\selectfont\ttfamily
          \ifSeriesDraft\SeriesDraftWord\SideStripeSep\fi%
          \SideStripeRepo\SideStripeSep results/#1\SideStripeSep
          \PublishedStripe\SideStripeSep last updated \BuildDateLong}}}}}

% CHANGE OF 2026-09-22, on the owner's word, to the one file all the papers of the series share.
%
% The stripe now carries TWO dates and names the folder as results/ rather than research/:
%
%   DRAFT . evilaliv3/vesuvius-challenge-pipeline . results/<folder> . <published> . last updated <date>
%
% \PublishedStripe comes from a published-date.tex beside the paper and is the day the folder
% entered the PUBLIC repository, written by hand once and never taken from the clock; a folder
% that is not published yet defines it as «not yet published». \BuildDateLong is the other
% quantity and keeps coming from SOURCE_DATE_EPOCH. A build refuses a publication date that is
% after the build date.
%
% EVERY PAPER THAT INPUTS THIS FILE MUST NOW ALSO INPUT A published-date.tex, or \PublishedStripe
% is undefined and the build stops. The three published papers of paper/00001 to 00003 are NOT
% rebuilt today: they are recompiled with their true publication dates at their next revision,
% which is the owner's decision of 2026-09-22 and is written in REVISION.md with its date.

%
% and in the BODY, immediately after \begin{document}, with the paper's own folder number:
%
%     \sidestripe{00001}
%
% The second line has to be in the body because the stripe is attached to the next page shipped
% out, which is the first page; the first line has to be in the preamble because it loads packages.
% The build of each paper fails if it carries neither.
%
% The date is \BuildDateLong, written by the build from SOURCE_DATE_EPOCH, the same value pdfTeX
% stamps into /CreationDate. It is never typed, so the front page and the metadata cannot come to
% say two different days. It is the date of the revision and not a reading of the clock, which is
% why the stripe labels it "last updated": a bare date beside a title is read as the day of
% writing, of submission or of the build, and a reader has no way to tell which was meant.

% THE DRAFT SWITCH OF THE SERIES. This one line is the whole control, for all three papers.
%
%     \SeriesDrafttrue    the papers are drafts and carry the mark
%     \SeriesDraftfalse   the papers are final and carry nothing
%
% Nothing else moves when it is thrown. The three papers reach this file through
% % The marks of the series: what a paper is, in the left margin of its first page, and, while the
% series is a draft, the word DRAFT across every page of it.
%
% One file for the three papers, so that a later change to the wording, the colour or the position
% lands in all three at once. It draws, rotated a quarter turn and read from the bottom, the state
% of the paper, the repository, the paper's own folder and the date of the revision it prints:
%
%     DRAFT . evilaliv3/vesuvius-challenge-pipeline . research/00001 . last updated 18 September 2026
%
% Why it exists. These are self published reports whose claims are pinned to a codebase that moves
% under them and to data that changes, and the method they argue for rests on when a thing was
% declared against when it was measured. A paper that hides its own date is inconsistent with the
% way it argues, and a reader who prints the front page should be able to see, without reading a
% line of it, what the thing is, what state it is in and when it was last revised.
%
% Usage, two lines per paper. In the PREAMBLE, after the class and before \begin{document}:
%
%     % The marks of the series: what a paper is, in the left margin of its first page, and, while the
% series is a draft, the word DRAFT across every page of it.
%
% One file for the three papers, so that a later change to the wording, the colour or the position
% lands in all three at once. It draws, rotated a quarter turn and read from the bottom, the state
% of the paper, the repository, the paper's own folder and the date of the revision it prints:
%
%     DRAFT . evilaliv3/vesuvius-challenge-pipeline . research/00001 . last updated 18 September 2026
%
% Why it exists. These are self published reports whose claims are pinned to a codebase that moves
% under them and to data that changes, and the method they argue for rests on when a thing was
% declared against when it was measured. A paper that hides its own date is inconsistent with the
% way it argues, and a reader who prints the front page should be able to see, without reading a
% line of it, what the thing is, what state it is in and when it was last revised.
%
% Usage, two lines per paper. In the PREAMBLE, after the class and before \begin{document}:
%
%     \input{sidestripe}
%
% and in the BODY, immediately after \begin{document}, with the paper's own folder number:
%
%     \sidestripe{00001}
%
% The second line has to be in the body because the stripe is attached to the next page shipped
% out, which is the first page; the first line has to be in the preamble because it loads packages.
% The build of each paper fails if it carries neither.
%
% The date is \BuildDateLong, written by the build from SOURCE_DATE_EPOCH, the same value pdfTeX
% stamps into /CreationDate. It is never typed, so the front page and the metadata cannot come to
% say two different days. It is the date of the revision and not a reading of the clock, which is
% why the stripe labels it "last updated": a bare date beside a title is read as the day of
% writing, of submission or of the build, and a reader has no way to tell which was meant.

% THE DRAFT SWITCH OF THE SERIES. This one line is the whole control, for all three papers.
%
%     \SeriesDrafttrue    the papers are drafts and carry the mark
%     \SeriesDraftfalse   the papers are final and carry nothing
%
% Nothing else moves when it is thrown. The three papers reach this file through
% \input{sidestripe}, and the build of each one refuses a paper that does not, so none of them can
% miss the switch and no paper of the series can be marked while another is not.
%
% Off, this file draws no watermark and loads no package it did not load before, and each paper
% rebuilds to the bytes it had before the mark existed. That was checked on 2026-09-18 by building
% all three with the switch off and comparing the sha256 with the published files, not assumed.
\newif\ifSeriesDraft
\SeriesDraftfalse

\usepackage{eso-pic}
\usepackage{xcolor}

% The build writes build-date.tex and every paper of the series inputs it before this file. If it
% is missing, the date is not declared anywhere, and the build stops here rather than reaching a
% page that quietly carries no date or, worse, one taken from the clock.
\ifdefined\BuildDateLong\else
  \PackageError{sidestripe}{The command BuildDateLong is not defined: input build-date first}
    {The build writes build-date.tex from SOURCE_DATE_EPOCH. Run the paper's build script.}
\fi

\newlength{\SideStripeX}\setlength{\SideStripeX}{0.42in}
\newlength{\SideStripeY}\setlength{\SideStripeY}{1.05in}
% The separator is a middle dot: the house rule forbids a dash as punctuation, and a comma would
% read as a list rather than as the three fields of an identifier.
\newcommand{\SideStripeSep}{\,\textperiodcentered\,}
\newcommand{\SideStripeRepo}{evilaliv3/vesuvius-challenge-pipeline}

% The watermark, drawn only while the switch is on.
%
% Why a mark on every page and not a banner on the first. A reader almost never meets these papers
% whole: a page is printed, a screenshot is pasted into a message, a section is forwarded. A mark
% that lives on the first page alone is gone at that moment, and the page that travels is the one
% that will be quoted back. A running head would reach every page too, but it belongs to the
% galley: in this two column class it would move the text block, change where the columns break
% and cost a page, and it would have to be removed from the galley again to get the final paper
% back. A shipout picture costs nothing of the kind.
%
% Why the FOREGROUND and not the background. It was drawn in the background until 2026-09-18, and
% there it was invisible wherever a figure sat, because every figure in these papers fills its own
% box with opaque white before it draws anything. A figure is exactly the part of a page somebody
% screenshots, so the mark was missing from the part most likely to travel alone. On top of
% everything it reaches the whole page, figures included.
%
% What keeps it from spoiling the page. It is drawn at \SeriesDraftOpacity of full opacity, so what
% lands on the paper is a tint and not a shape: over white it is the lightest grey a PDF can carry,
% and over black text it changes nothing at all, because a tint of black over black is black. The
% text underneath is therefore not dimmed in any way a reader can measure. Nothing is added to the
% galley either, so the page count, the line breaks and the column breaks are exactly those of the
% final paper, and the gate on overfull boxes sees nothing new.
%
% The colour is black and not a grey, because at this opacity the colour is the only thing left to
% spend: the tint over white is 1 minus the opacity of the way from white to the ink, so black is
% the darkest mark two per cent can buy. A light grey at two per cent would be indistinguishable
% from the paper.
%
% \resizebox scales a box set at the normal size instead of asking for a font at 80pt, which the
% Computer Modern shapes of this class do not have: LaTeX would substitute the nearest design size
% and print a mark several times smaller than the one asked for, with a font warning the gates do
% not read. Scaling the box cannot fail that way, and it is a linear transform in the PDF, so the
% result is the same bytes on every build.
\ifSeriesDraft
  % transparent.sty writes the constant alpha into an ExtGState, which is what pdfTeX needs to put
  % ink on top of ink without hiding it. It is loaded only here, so a paper built with the switch
  % off carries no transparency machinery and no extra resource at all.
  \usepackage{transparent}
  \newcommand{\SeriesDraftWord}{DRAFT}
  \newcommand{\SeriesDraftOpacity}{0.10}
  \AddToShipoutPictureFG{%
    \AtPageCenter{%
      \makebox(0,0){%
        \rotatebox{45}{%
          \transparent{\SeriesDraftOpacity}%
          \resizebox{0.68\paperwidth}{!}{%
            \color{black}\sffamily\bfseries\SeriesDraftWord}}}}}
\fi

% THE CASE OF THE SERIES: the headings and the table captions read as the rest of the article.
%
% IEEEtran sets two things in small capitals, and a reader sees them as ALL CAPITALS shouted in the
% middle of a page whose every other line is ordinary sentence case:
%
%   the table captions, IEEEtran.cls line 2778, which puts the caption text through \scshape;
%   the section headings, IEEEtran.cls line 5473, whose heading font ends in \scshape.
%
% Those are the two the responsible person asked for on 2026-09-18, and they are the only two that
% are live in this class option. The other \scshape sites in the class, lines 2737, 5468, 5501 and
% 5790, belong to the conference and computer society options, which these papers do not use, and
% the figure captions never had it: figure captions were already sentence case while table captions
% were not, so this change also makes the two agree, which they did not before.
%
% What is deliberately NOT changed: the drop cap at the start of each section, IEEEtran.cls lines
% 5835 and 5841, which sets the large initial and the rest of its first word through
% \MakeUppercase. A large initial followed by a word in capitals is a typographic convention and
% not shouting, and it is a bigger aesthetic move than was asked for. One line in this file takes
% it off if it is ever wanted: \renewcommand{\IEEEPARstartWORDCAPSTYLE}{\relax}.
%
% The class file is NOT edited. It is a third party class, vendored in three trees, and a modified
% copy under the same name is what makes a build unreproducible for somebody else. etoolbox patches
% the two macros from the preamble instead, so one change here reaches all three papers exactly as
% the draft switch does.
%
% Both patches fail the build if they do not apply. That matters more than it looks: if the class
% is ever updated and the macro no longer matches, a silent failure would bring the capitals back
% and nobody would know which build changed them.
\usepackage{etoolbox}
\makeatletter
\patchcmd{\@makecaption}{\scshape}{}{}{%
  \PackageError{sidestripe}{Cannot take the small capitals out of the table captions}
    {IEEEtran's \string\@makecaption no longer contains \string\scshape where it did at line
     2778. Read the class and fix the patch rather than removing it.}}
\patchcmd{\section}{\scshape}{}{}{%
  \PackageError{sidestripe}{Cannot take the small capitals out of the section headings}
    {IEEEtran's \string\section no longer contains \string\scshape where it did at line 5473.
     Read the class and fix the patch rather than removing it.}}
\makeatother

% Mechanics of the stripe. eso-pic's \AddToShipoutPictureBG* draws in the background layer of the
% current page only, which is what keeps the body untouched: the text block is not moved by a
% single point, because nothing is added to the galley. \rotatebox comes from graphicx, which the
% papers already load for their figures, and the grey is xcolor. tikz with `remember picture,
% overlay' would do the same thing at the price of a second compilation pass and a much larger
% dependency, and a raw \rotatebox in a zero size box inside the body would have to be placed by
% hand relative to the text and would move with it.
%
% Geometry. The stripe sits \SideStripeX from the left edge of the paper and starts \SideStripeY
% from the bottom, reading upwards. Those are far enough inside the paper that a printer's trim
% does not eat them and far enough from the text block that they cannot collide with it; both are
% named here so that the three papers cannot drift apart.
% If the paper that inputs this file has no published-date.tex, say so in words rather than
% letting TeX report an undefined control sequence from inside a rotated box. The three papers of
% paper/00001 to 00003 are in exactly that state on 2026-09-22 and are not rebuilt today; whoever
% builds one next is told here what to add and why, instead of debugging it.
\ifdefined\PublishedStripe\else
  \errmessage{sidestripe.tex: this paper inputs the stripe but defines no \string\PublishedStripe.
  Add a published-date.tex beside the paper and \string\input{published-date}. It holds the day
  the folder entered the PUBLIC repository, written by hand and never from the clock; a folder not
  yet published defines \string\PublishedStripe as: not yet published. See REVISION.md,
  2026-09-22}
\fi

\newcommand{\sidestripe}[1]{%
  \AddToShipoutPictureBG*{%
    \AtPageLowerLeft{%
      \put(\LenToUnit{\SideStripeX},\LenToUnit{\SideStripeY}){%
        \rotatebox{90}{%
          \color[gray]{0.55}%
          \fontsize{7}{8}\selectfont\ttfamily
          \ifSeriesDraft\SeriesDraftWord\SideStripeSep\fi%
          \SideStripeRepo\SideStripeSep results/#1\SideStripeSep
          \PublishedStripe\SideStripeSep last updated \BuildDateLong}}}}}

% CHANGE OF 2026-09-22, on the owner's word, to the one file all the papers of the series share.
%
% The stripe now carries TWO dates and names the folder as results/ rather than research/:
%
%   DRAFT . evilaliv3/vesuvius-challenge-pipeline . results/<folder> . <published> . last updated <date>
%
% \PublishedStripe comes from a published-date.tex beside the paper and is the day the folder
% entered the PUBLIC repository, written by hand once and never taken from the clock; a folder
% that is not published yet defines it as «not yet published». \BuildDateLong is the other
% quantity and keeps coming from SOURCE_DATE_EPOCH. A build refuses a publication date that is
% after the build date.
%
% EVERY PAPER THAT INPUTS THIS FILE MUST NOW ALSO INPUT A published-date.tex, or \PublishedStripe
% is undefined and the build stops. The three published papers of paper/00001 to 00003 are NOT
% rebuilt today: they are recompiled with their true publication dates at their next revision,
% which is the owner's decision of 2026-09-22 and is written in REVISION.md with its date.

%
% and in the BODY, immediately after \begin{document}, with the paper's own folder number:
%
%     \sidestripe{00001}
%
% The second line has to be in the body because the stripe is attached to the next page shipped
% out, which is the first page; the first line has to be in the preamble because it loads packages.
% The build of each paper fails if it carries neither.
%
% The date is \BuildDateLong, written by the build from SOURCE_DATE_EPOCH, the same value pdfTeX
% stamps into /CreationDate. It is never typed, so the front page and the metadata cannot come to
% say two different days. It is the date of the revision and not a reading of the clock, which is
% why the stripe labels it "last updated": a bare date beside a title is read as the day of
% writing, of submission or of the build, and a reader has no way to tell which was meant.

% THE DRAFT SWITCH OF THE SERIES. This one line is the whole control, for all three papers.
%
%     \SeriesDrafttrue    the papers are drafts and carry the mark
%     \SeriesDraftfalse   the papers are final and carry nothing
%
% Nothing else moves when it is thrown. The three papers reach this file through
% % The marks of the series: what a paper is, in the left margin of its first page, and, while the
% series is a draft, the word DRAFT across every page of it.
%
% One file for the three papers, so that a later change to the wording, the colour or the position
% lands in all three at once. It draws, rotated a quarter turn and read from the bottom, the state
% of the paper, the repository, the paper's own folder and the date of the revision it prints:
%
%     DRAFT . evilaliv3/vesuvius-challenge-pipeline . research/00001 . last updated 18 September 2026
%
% Why it exists. These are self published reports whose claims are pinned to a codebase that moves
% under them and to data that changes, and the method they argue for rests on when a thing was
% declared against when it was measured. A paper that hides its own date is inconsistent with the
% way it argues, and a reader who prints the front page should be able to see, without reading a
% line of it, what the thing is, what state it is in and when it was last revised.
%
% Usage, two lines per paper. In the PREAMBLE, after the class and before \begin{document}:
%
%     \input{sidestripe}
%
% and in the BODY, immediately after \begin{document}, with the paper's own folder number:
%
%     \sidestripe{00001}
%
% The second line has to be in the body because the stripe is attached to the next page shipped
% out, which is the first page; the first line has to be in the preamble because it loads packages.
% The build of each paper fails if it carries neither.
%
% The date is \BuildDateLong, written by the build from SOURCE_DATE_EPOCH, the same value pdfTeX
% stamps into /CreationDate. It is never typed, so the front page and the metadata cannot come to
% say two different days. It is the date of the revision and not a reading of the clock, which is
% why the stripe labels it "last updated": a bare date beside a title is read as the day of
% writing, of submission or of the build, and a reader has no way to tell which was meant.

% THE DRAFT SWITCH OF THE SERIES. This one line is the whole control, for all three papers.
%
%     \SeriesDrafttrue    the papers are drafts and carry the mark
%     \SeriesDraftfalse   the papers are final and carry nothing
%
% Nothing else moves when it is thrown. The three papers reach this file through
% \input{sidestripe}, and the build of each one refuses a paper that does not, so none of them can
% miss the switch and no paper of the series can be marked while another is not.
%
% Off, this file draws no watermark and loads no package it did not load before, and each paper
% rebuilds to the bytes it had before the mark existed. That was checked on 2026-09-18 by building
% all three with the switch off and comparing the sha256 with the published files, not assumed.
\newif\ifSeriesDraft
\SeriesDraftfalse

\usepackage{eso-pic}
\usepackage{xcolor}

% The build writes build-date.tex and every paper of the series inputs it before this file. If it
% is missing, the date is not declared anywhere, and the build stops here rather than reaching a
% page that quietly carries no date or, worse, one taken from the clock.
\ifdefined\BuildDateLong\else
  \PackageError{sidestripe}{The command BuildDateLong is not defined: input build-date first}
    {The build writes build-date.tex from SOURCE_DATE_EPOCH. Run the paper's build script.}
\fi

\newlength{\SideStripeX}\setlength{\SideStripeX}{0.42in}
\newlength{\SideStripeY}\setlength{\SideStripeY}{1.05in}
% The separator is a middle dot: the house rule forbids a dash as punctuation, and a comma would
% read as a list rather than as the three fields of an identifier.
\newcommand{\SideStripeSep}{\,\textperiodcentered\,}
\newcommand{\SideStripeRepo}{evilaliv3/vesuvius-challenge-pipeline}

% The watermark, drawn only while the switch is on.
%
% Why a mark on every page and not a banner on the first. A reader almost never meets these papers
% whole: a page is printed, a screenshot is pasted into a message, a section is forwarded. A mark
% that lives on the first page alone is gone at that moment, and the page that travels is the one
% that will be quoted back. A running head would reach every page too, but it belongs to the
% galley: in this two column class it would move the text block, change where the columns break
% and cost a page, and it would have to be removed from the galley again to get the final paper
% back. A shipout picture costs nothing of the kind.
%
% Why the FOREGROUND and not the background. It was drawn in the background until 2026-09-18, and
% there it was invisible wherever a figure sat, because every figure in these papers fills its own
% box with opaque white before it draws anything. A figure is exactly the part of a page somebody
% screenshots, so the mark was missing from the part most likely to travel alone. On top of
% everything it reaches the whole page, figures included.
%
% What keeps it from spoiling the page. It is drawn at \SeriesDraftOpacity of full opacity, so what
% lands on the paper is a tint and not a shape: over white it is the lightest grey a PDF can carry,
% and over black text it changes nothing at all, because a tint of black over black is black. The
% text underneath is therefore not dimmed in any way a reader can measure. Nothing is added to the
% galley either, so the page count, the line breaks and the column breaks are exactly those of the
% final paper, and the gate on overfull boxes sees nothing new.
%
% The colour is black and not a grey, because at this opacity the colour is the only thing left to
% spend: the tint over white is 1 minus the opacity of the way from white to the ink, so black is
% the darkest mark two per cent can buy. A light grey at two per cent would be indistinguishable
% from the paper.
%
% \resizebox scales a box set at the normal size instead of asking for a font at 80pt, which the
% Computer Modern shapes of this class do not have: LaTeX would substitute the nearest design size
% and print a mark several times smaller than the one asked for, with a font warning the gates do
% not read. Scaling the box cannot fail that way, and it is a linear transform in the PDF, so the
% result is the same bytes on every build.
\ifSeriesDraft
  % transparent.sty writes the constant alpha into an ExtGState, which is what pdfTeX needs to put
  % ink on top of ink without hiding it. It is loaded only here, so a paper built with the switch
  % off carries no transparency machinery and no extra resource at all.
  \usepackage{transparent}
  \newcommand{\SeriesDraftWord}{DRAFT}
  \newcommand{\SeriesDraftOpacity}{0.10}
  \AddToShipoutPictureFG{%
    \AtPageCenter{%
      \makebox(0,0){%
        \rotatebox{45}{%
          \transparent{\SeriesDraftOpacity}%
          \resizebox{0.68\paperwidth}{!}{%
            \color{black}\sffamily\bfseries\SeriesDraftWord}}}}}
\fi

% THE CASE OF THE SERIES: the headings and the table captions read as the rest of the article.
%
% IEEEtran sets two things in small capitals, and a reader sees them as ALL CAPITALS shouted in the
% middle of a page whose every other line is ordinary sentence case:
%
%   the table captions, IEEEtran.cls line 2778, which puts the caption text through \scshape;
%   the section headings, IEEEtran.cls line 5473, whose heading font ends in \scshape.
%
% Those are the two the responsible person asked for on 2026-09-18, and they are the only two that
% are live in this class option. The other \scshape sites in the class, lines 2737, 5468, 5501 and
% 5790, belong to the conference and computer society options, which these papers do not use, and
% the figure captions never had it: figure captions were already sentence case while table captions
% were not, so this change also makes the two agree, which they did not before.
%
% What is deliberately NOT changed: the drop cap at the start of each section, IEEEtran.cls lines
% 5835 and 5841, which sets the large initial and the rest of its first word through
% \MakeUppercase. A large initial followed by a word in capitals is a typographic convention and
% not shouting, and it is a bigger aesthetic move than was asked for. One line in this file takes
% it off if it is ever wanted: \renewcommand{\IEEEPARstartWORDCAPSTYLE}{\relax}.
%
% The class file is NOT edited. It is a third party class, vendored in three trees, and a modified
% copy under the same name is what makes a build unreproducible for somebody else. etoolbox patches
% the two macros from the preamble instead, so one change here reaches all three papers exactly as
% the draft switch does.
%
% Both patches fail the build if they do not apply. That matters more than it looks: if the class
% is ever updated and the macro no longer matches, a silent failure would bring the capitals back
% and nobody would know which build changed them.
\usepackage{etoolbox}
\makeatletter
\patchcmd{\@makecaption}{\scshape}{}{}{%
  \PackageError{sidestripe}{Cannot take the small capitals out of the table captions}
    {IEEEtran's \string\@makecaption no longer contains \string\scshape where it did at line
     2778. Read the class and fix the patch rather than removing it.}}
\patchcmd{\section}{\scshape}{}{}{%
  \PackageError{sidestripe}{Cannot take the small capitals out of the section headings}
    {IEEEtran's \string\section no longer contains \string\scshape where it did at line 5473.
     Read the class and fix the patch rather than removing it.}}
\makeatother

% Mechanics of the stripe. eso-pic's \AddToShipoutPictureBG* draws in the background layer of the
% current page only, which is what keeps the body untouched: the text block is not moved by a
% single point, because nothing is added to the galley. \rotatebox comes from graphicx, which the
% papers already load for their figures, and the grey is xcolor. tikz with `remember picture,
% overlay' would do the same thing at the price of a second compilation pass and a much larger
% dependency, and a raw \rotatebox in a zero size box inside the body would have to be placed by
% hand relative to the text and would move with it.
%
% Geometry. The stripe sits \SideStripeX from the left edge of the paper and starts \SideStripeY
% from the bottom, reading upwards. Those are far enough inside the paper that a printer's trim
% does not eat them and far enough from the text block that they cannot collide with it; both are
% named here so that the three papers cannot drift apart.
% If the paper that inputs this file has no published-date.tex, say so in words rather than
% letting TeX report an undefined control sequence from inside a rotated box. The three papers of
% paper/00001 to 00003 are in exactly that state on 2026-09-22 and are not rebuilt today; whoever
% builds one next is told here what to add and why, instead of debugging it.
\ifdefined\PublishedStripe\else
  \errmessage{sidestripe.tex: this paper inputs the stripe but defines no \string\PublishedStripe.
  Add a published-date.tex beside the paper and \string\input{published-date}. It holds the day
  the folder entered the PUBLIC repository, written by hand and never from the clock; a folder not
  yet published defines \string\PublishedStripe as: not yet published. See REVISION.md,
  2026-09-22}
\fi

\newcommand{\sidestripe}[1]{%
  \AddToShipoutPictureBG*{%
    \AtPageLowerLeft{%
      \put(\LenToUnit{\SideStripeX},\LenToUnit{\SideStripeY}){%
        \rotatebox{90}{%
          \color[gray]{0.55}%
          \fontsize{7}{8}\selectfont\ttfamily
          \ifSeriesDraft\SeriesDraftWord\SideStripeSep\fi%
          \SideStripeRepo\SideStripeSep results/#1\SideStripeSep
          \PublishedStripe\SideStripeSep last updated \BuildDateLong}}}}}

% CHANGE OF 2026-09-22, on the owner's word, to the one file all the papers of the series share.
%
% The stripe now carries TWO dates and names the folder as results/ rather than research/:
%
%   DRAFT . evilaliv3/vesuvius-challenge-pipeline . results/<folder> . <published> . last updated <date>
%
% \PublishedStripe comes from a published-date.tex beside the paper and is the day the folder
% entered the PUBLIC repository, written by hand once and never taken from the clock; a folder
% that is not published yet defines it as «not yet published». \BuildDateLong is the other
% quantity and keeps coming from SOURCE_DATE_EPOCH. A build refuses a publication date that is
% after the build date.
%
% EVERY PAPER THAT INPUTS THIS FILE MUST NOW ALSO INPUT A published-date.tex, or \PublishedStripe
% is undefined and the build stops. The three published papers of paper/00001 to 00003 are NOT
% rebuilt today: they are recompiled with their true publication dates at their next revision,
% which is the owner's decision of 2026-09-22 and is written in REVISION.md with its date.
, and the build of each one refuses a paper that does not, so none of them can
% miss the switch and no paper of the series can be marked while another is not.
%
% Off, this file draws no watermark and loads no package it did not load before, and each paper
% rebuilds to the bytes it had before the mark existed. That was checked on 2026-09-18 by building
% all three with the switch off and comparing the sha256 with the published files, not assumed.
\newif\ifSeriesDraft
\SeriesDraftfalse

\usepackage{eso-pic}
\usepackage{xcolor}

% The build writes build-date.tex and every paper of the series inputs it before this file. If it
% is missing, the date is not declared anywhere, and the build stops here rather than reaching a
% page that quietly carries no date or, worse, one taken from the clock.
\ifdefined\BuildDateLong\else
  \PackageError{sidestripe}{The command BuildDateLong is not defined: input build-date first}
    {The build writes build-date.tex from SOURCE_DATE_EPOCH. Run the paper's build script.}
\fi

\newlength{\SideStripeX}\setlength{\SideStripeX}{0.42in}
\newlength{\SideStripeY}\setlength{\SideStripeY}{1.05in}
% The separator is a middle dot: the house rule forbids a dash as punctuation, and a comma would
% read as a list rather than as the three fields of an identifier.
\newcommand{\SideStripeSep}{\,\textperiodcentered\,}
\newcommand{\SideStripeRepo}{evilaliv3/vesuvius-challenge-pipeline}

% The watermark, drawn only while the switch is on.
%
% Why a mark on every page and not a banner on the first. A reader almost never meets these papers
% whole: a page is printed, a screenshot is pasted into a message, a section is forwarded. A mark
% that lives on the first page alone is gone at that moment, and the page that travels is the one
% that will be quoted back. A running head would reach every page too, but it belongs to the
% galley: in this two column class it would move the text block, change where the columns break
% and cost a page, and it would have to be removed from the galley again to get the final paper
% back. A shipout picture costs nothing of the kind.
%
% Why the FOREGROUND and not the background. It was drawn in the background until 2026-09-18, and
% there it was invisible wherever a figure sat, because every figure in these papers fills its own
% box with opaque white before it draws anything. A figure is exactly the part of a page somebody
% screenshots, so the mark was missing from the part most likely to travel alone. On top of
% everything it reaches the whole page, figures included.
%
% What keeps it from spoiling the page. It is drawn at \SeriesDraftOpacity of full opacity, so what
% lands on the paper is a tint and not a shape: over white it is the lightest grey a PDF can carry,
% and over black text it changes nothing at all, because a tint of black over black is black. The
% text underneath is therefore not dimmed in any way a reader can measure. Nothing is added to the
% galley either, so the page count, the line breaks and the column breaks are exactly those of the
% final paper, and the gate on overfull boxes sees nothing new.
%
% The colour is black and not a grey, because at this opacity the colour is the only thing left to
% spend: the tint over white is 1 minus the opacity of the way from white to the ink, so black is
% the darkest mark two per cent can buy. A light grey at two per cent would be indistinguishable
% from the paper.
%
% \resizebox scales a box set at the normal size instead of asking for a font at 80pt, which the
% Computer Modern shapes of this class do not have: LaTeX would substitute the nearest design size
% and print a mark several times smaller than the one asked for, with a font warning the gates do
% not read. Scaling the box cannot fail that way, and it is a linear transform in the PDF, so the
% result is the same bytes on every build.
\ifSeriesDraft
  % transparent.sty writes the constant alpha into an ExtGState, which is what pdfTeX needs to put
  % ink on top of ink without hiding it. It is loaded only here, so a paper built with the switch
  % off carries no transparency machinery and no extra resource at all.
  \usepackage{transparent}
  \newcommand{\SeriesDraftWord}{DRAFT}
  \newcommand{\SeriesDraftOpacity}{0.10}
  \AddToShipoutPictureFG{%
    \AtPageCenter{%
      \makebox(0,0){%
        \rotatebox{45}{%
          \transparent{\SeriesDraftOpacity}%
          \resizebox{0.68\paperwidth}{!}{%
            \color{black}\sffamily\bfseries\SeriesDraftWord}}}}}
\fi

% THE CASE OF THE SERIES: the headings and the table captions read as the rest of the article.
%
% IEEEtran sets two things in small capitals, and a reader sees them as ALL CAPITALS shouted in the
% middle of a page whose every other line is ordinary sentence case:
%
%   the table captions, IEEEtran.cls line 2778, which puts the caption text through \scshape;
%   the section headings, IEEEtran.cls line 5473, whose heading font ends in \scshape.
%
% Those are the two the responsible person asked for on 2026-09-18, and they are the only two that
% are live in this class option. The other \scshape sites in the class, lines 2737, 5468, 5501 and
% 5790, belong to the conference and computer society options, which these papers do not use, and
% the figure captions never had it: figure captions were already sentence case while table captions
% were not, so this change also makes the two agree, which they did not before.
%
% What is deliberately NOT changed: the drop cap at the start of each section, IEEEtran.cls lines
% 5835 and 5841, which sets the large initial and the rest of its first word through
% \MakeUppercase. A large initial followed by a word in capitals is a typographic convention and
% not shouting, and it is a bigger aesthetic move than was asked for. One line in this file takes
% it off if it is ever wanted: \renewcommand{\IEEEPARstartWORDCAPSTYLE}{\relax}.
%
% The class file is NOT edited. It is a third party class, vendored in three trees, and a modified
% copy under the same name is what makes a build unreproducible for somebody else. etoolbox patches
% the two macros from the preamble instead, so one change here reaches all three papers exactly as
% the draft switch does.
%
% Both patches fail the build if they do not apply. That matters more than it looks: if the class
% is ever updated and the macro no longer matches, a silent failure would bring the capitals back
% and nobody would know which build changed them.
\usepackage{etoolbox}
\makeatletter
\patchcmd{\@makecaption}{\scshape}{}{}{%
  \PackageError{sidestripe}{Cannot take the small capitals out of the table captions}
    {IEEEtran's \string\@makecaption no longer contains \string\scshape where it did at line
     2778. Read the class and fix the patch rather than removing it.}}
\patchcmd{\section}{\scshape}{}{}{%
  \PackageError{sidestripe}{Cannot take the small capitals out of the section headings}
    {IEEEtran's \string\section no longer contains \string\scshape where it did at line 5473.
     Read the class and fix the patch rather than removing it.}}
\makeatother

% Mechanics of the stripe. eso-pic's \AddToShipoutPictureBG* draws in the background layer of the
% current page only, which is what keeps the body untouched: the text block is not moved by a
% single point, because nothing is added to the galley. \rotatebox comes from graphicx, which the
% papers already load for their figures, and the grey is xcolor. tikz with `remember picture,
% overlay' would do the same thing at the price of a second compilation pass and a much larger
% dependency, and a raw \rotatebox in a zero size box inside the body would have to be placed by
% hand relative to the text and would move with it.
%
% Geometry. The stripe sits \SideStripeX from the left edge of the paper and starts \SideStripeY
% from the bottom, reading upwards. Those are far enough inside the paper that a printer's trim
% does not eat them and far enough from the text block that they cannot collide with it; both are
% named here so that the three papers cannot drift apart.
% If the paper that inputs this file has no published-date.tex, say so in words rather than
% letting TeX report an undefined control sequence from inside a rotated box. The three papers of
% paper/00001 to 00003 are in exactly that state on 2026-09-22 and are not rebuilt today; whoever
% builds one next is told here what to add and why, instead of debugging it.
\ifdefined\PublishedStripe\else
  \errmessage{sidestripe.tex: this paper inputs the stripe but defines no \string\PublishedStripe.
  Add a published-date.tex beside the paper and \string% The day this folder entered the PUBLIC repository, written here by hand once and never taken
% from the clock: a build that read the date from the machine would move it every time and a
% front page would claim a day it was not published on. The last updated date beside it in the
% stripe is a different quantity and comes from SOURCE_DATE_EPOCH.
%
% This folder is NOT in the public repository yet: on 2026-09-22 the only folder under
% results/ of evilaliv3/vesuvius-challenge-pipeline is e89eb7ad-umbilicus-estimator. So the
% stripe carries «not yet published» rather than a date, and when the folder is published the
% two lines below are replaced with the day it was, in this shape:
%
%   \newcommand{\PublishedDateIso}{2026-09-22}
%   \newcommand{\PublishedStripe}{published 22 September 2026}
%
% build.sh reads \PublishedDateIso and refuses a date in the future.
\newcommand{\PublishedDateIso}{}
\newcommand{\PublishedStripe}{not yet published}
. It holds the day
  the folder entered the PUBLIC repository, written by hand and never from the clock; a folder not
  yet published defines \string\PublishedStripe as: not yet published. See REVISION.md,
  2026-09-22}
\fi

\newcommand{\sidestripe}[1]{%
  \AddToShipoutPictureBG*{%
    \AtPageLowerLeft{%
      \put(\LenToUnit{\SideStripeX},\LenToUnit{\SideStripeY}){%
        \rotatebox{90}{%
          \color[gray]{0.55}%
          \fontsize{7}{8}\selectfont\ttfamily
          \ifSeriesDraft\SeriesDraftWord\SideStripeSep\fi%
          \SideStripeRepo\SideStripeSep results/#1\SideStripeSep
          \PublishedStripe\SideStripeSep last updated \BuildDateLong}}}}}

% CHANGE OF 2026-09-22, on the owner's word, to the one file all the papers of the series share.
%
% The stripe now carries TWO dates and names the folder as results/ rather than research/:
%
%   DRAFT . evilaliv3/vesuvius-challenge-pipeline . results/<folder> . <published> . last updated <date>
%
% \PublishedStripe comes from a published-date.tex beside the paper and is the day the folder
% entered the PUBLIC repository, written by hand once and never taken from the clock; a folder
% that is not published yet defines it as «not yet published». \BuildDateLong is the other
% quantity and keeps coming from SOURCE_DATE_EPOCH. A build refuses a publication date that is
% after the build date.
%
% EVERY PAPER THAT INPUTS THIS FILE MUST NOW ALSO INPUT A published-date.tex, or \PublishedStripe
% is undefined and the build stops. The three published papers of paper/00001 to 00003 are NOT
% rebuilt today: they are recompiled with their true publication dates at their next revision,
% which is the owner's decision of 2026-09-22 and is written in REVISION.md with its date.
, and the build of each one refuses a paper that does not, so none of them can
% miss the switch and no paper of the series can be marked while another is not.
%
% Off, this file draws no watermark and loads no package it did not load before, and each paper
% rebuilds to the bytes it had before the mark existed. That was checked on 2026-09-18 by building
% all three with the switch off and comparing the sha256 with the published files, not assumed.
\newif\ifSeriesDraft
\SeriesDraftfalse

\usepackage{eso-pic}
\usepackage{xcolor}

% The build writes build-date.tex and every paper of the series inputs it before this file. If it
% is missing, the date is not declared anywhere, and the build stops here rather than reaching a
% page that quietly carries no date or, worse, one taken from the clock.
\ifdefined\BuildDateLong\else
  \PackageError{sidestripe}{The command BuildDateLong is not defined: input build-date first}
    {The build writes build-date.tex from SOURCE_DATE_EPOCH. Run the paper's build script.}
\fi

\newlength{\SideStripeX}\setlength{\SideStripeX}{0.42in}
\newlength{\SideStripeY}\setlength{\SideStripeY}{1.05in}
% The separator is a middle dot: the house rule forbids a dash as punctuation, and a comma would
% read as a list rather than as the three fields of an identifier.
\newcommand{\SideStripeSep}{\,\textperiodcentered\,}
\newcommand{\SideStripeRepo}{evilaliv3/vesuvius-challenge-pipeline}

% The watermark, drawn only while the switch is on.
%
% Why a mark on every page and not a banner on the first. A reader almost never meets these papers
% whole: a page is printed, a screenshot is pasted into a message, a section is forwarded. A mark
% that lives on the first page alone is gone at that moment, and the page that travels is the one
% that will be quoted back. A running head would reach every page too, but it belongs to the
% galley: in this two column class it would move the text block, change where the columns break
% and cost a page, and it would have to be removed from the galley again to get the final paper
% back. A shipout picture costs nothing of the kind.
%
% Why the FOREGROUND and not the background. It was drawn in the background until 2026-09-18, and
% there it was invisible wherever a figure sat, because every figure in these papers fills its own
% box with opaque white before it draws anything. A figure is exactly the part of a page somebody
% screenshots, so the mark was missing from the part most likely to travel alone. On top of
% everything it reaches the whole page, figures included.
%
% What keeps it from spoiling the page. It is drawn at \SeriesDraftOpacity of full opacity, so what
% lands on the paper is a tint and not a shape: over white it is the lightest grey a PDF can carry,
% and over black text it changes nothing at all, because a tint of black over black is black. The
% text underneath is therefore not dimmed in any way a reader can measure. Nothing is added to the
% galley either, so the page count, the line breaks and the column breaks are exactly those of the
% final paper, and the gate on overfull boxes sees nothing new.
%
% The colour is black and not a grey, because at this opacity the colour is the only thing left to
% spend: the tint over white is 1 minus the opacity of the way from white to the ink, so black is
% the darkest mark two per cent can buy. A light grey at two per cent would be indistinguishable
% from the paper.
%
% \resizebox scales a box set at the normal size instead of asking for a font at 80pt, which the
% Computer Modern shapes of this class do not have: LaTeX would substitute the nearest design size
% and print a mark several times smaller than the one asked for, with a font warning the gates do
% not read. Scaling the box cannot fail that way, and it is a linear transform in the PDF, so the
% result is the same bytes on every build.
\ifSeriesDraft
  % transparent.sty writes the constant alpha into an ExtGState, which is what pdfTeX needs to put
  % ink on top of ink without hiding it. It is loaded only here, so a paper built with the switch
  % off carries no transparency machinery and no extra resource at all.
  \usepackage{transparent}
  \newcommand{\SeriesDraftWord}{DRAFT}
  \newcommand{\SeriesDraftOpacity}{0.10}
  \AddToShipoutPictureFG{%
    \AtPageCenter{%
      \makebox(0,0){%
        \rotatebox{45}{%
          \transparent{\SeriesDraftOpacity}%
          \resizebox{0.68\paperwidth}{!}{%
            \color{black}\sffamily\bfseries\SeriesDraftWord}}}}}
\fi

% THE CASE OF THE SERIES: the headings and the table captions read as the rest of the article.
%
% IEEEtran sets two things in small capitals, and a reader sees them as ALL CAPITALS shouted in the
% middle of a page whose every other line is ordinary sentence case:
%
%   the table captions, IEEEtran.cls line 2778, which puts the caption text through \scshape;
%   the section headings, IEEEtran.cls line 5473, whose heading font ends in \scshape.
%
% Those are the two the responsible person asked for on 2026-09-18, and they are the only two that
% are live in this class option. The other \scshape sites in the class, lines 2737, 5468, 5501 and
% 5790, belong to the conference and computer society options, which these papers do not use, and
% the figure captions never had it: figure captions were already sentence case while table captions
% were not, so this change also makes the two agree, which they did not before.
%
% What is deliberately NOT changed: the drop cap at the start of each section, IEEEtran.cls lines
% 5835 and 5841, which sets the large initial and the rest of its first word through
% \MakeUppercase. A large initial followed by a word in capitals is a typographic convention and
% not shouting, and it is a bigger aesthetic move than was asked for. One line in this file takes
% it off if it is ever wanted: \renewcommand{\IEEEPARstartWORDCAPSTYLE}{\relax}.
%
% The class file is NOT edited. It is a third party class, vendored in three trees, and a modified
% copy under the same name is what makes a build unreproducible for somebody else. etoolbox patches
% the two macros from the preamble instead, so one change here reaches all three papers exactly as
% the draft switch does.
%
% Both patches fail the build if they do not apply. That matters more than it looks: if the class
% is ever updated and the macro no longer matches, a silent failure would bring the capitals back
% and nobody would know which build changed them.
\usepackage{etoolbox}
\makeatletter
\patchcmd{\@makecaption}{\scshape}{}{}{%
  \PackageError{sidestripe}{Cannot take the small capitals out of the table captions}
    {IEEEtran's \string\@makecaption no longer contains \string\scshape where it did at line
     2778. Read the class and fix the patch rather than removing it.}}
\patchcmd{\section}{\scshape}{}{}{%
  \PackageError{sidestripe}{Cannot take the small capitals out of the section headings}
    {IEEEtran's \string\section no longer contains \string\scshape where it did at line 5473.
     Read the class and fix the patch rather than removing it.}}
\makeatother

% Mechanics of the stripe. eso-pic's \AddToShipoutPictureBG* draws in the background layer of the
% current page only, which is what keeps the body untouched: the text block is not moved by a
% single point, because nothing is added to the galley. \rotatebox comes from graphicx, which the
% papers already load for their figures, and the grey is xcolor. tikz with `remember picture,
% overlay' would do the same thing at the price of a second compilation pass and a much larger
% dependency, and a raw \rotatebox in a zero size box inside the body would have to be placed by
% hand relative to the text and would move with it.
%
% Geometry. The stripe sits \SideStripeX from the left edge of the paper and starts \SideStripeY
% from the bottom, reading upwards. Those are far enough inside the paper that a printer's trim
% does not eat them and far enough from the text block that they cannot collide with it; both are
% named here so that the three papers cannot drift apart.
% If the paper that inputs this file has no published-date.tex, say so in words rather than
% letting TeX report an undefined control sequence from inside a rotated box. The three papers of
% paper/00001 to 00003 are in exactly that state on 2026-09-22 and are not rebuilt today; whoever
% builds one next is told here what to add and why, instead of debugging it.
\ifdefined\PublishedStripe\else
  \errmessage{sidestripe.tex: this paper inputs the stripe but defines no \string\PublishedStripe.
  Add a published-date.tex beside the paper and \string% The day this folder entered the PUBLIC repository, written here by hand once and never taken
% from the clock: a build that read the date from the machine would move it every time and a
% front page would claim a day it was not published on. The last updated date beside it in the
% stripe is a different quantity and comes from SOURCE_DATE_EPOCH.
%
% This folder is NOT in the public repository yet: on 2026-09-22 the only folder under
% results/ of evilaliv3/vesuvius-challenge-pipeline is e89eb7ad-umbilicus-estimator. So the
% stripe carries «not yet published» rather than a date, and when the folder is published the
% two lines below are replaced with the day it was, in this shape:
%
%   \newcommand{\PublishedDateIso}{2026-09-22}
%   \newcommand{\PublishedStripe}{published 22 September 2026}
%
% build.sh reads \PublishedDateIso and refuses a date in the future.
\newcommand{\PublishedDateIso}{}
\newcommand{\PublishedStripe}{not yet published}
. It holds the day
  the folder entered the PUBLIC repository, written by hand and never from the clock; a folder not
  yet published defines \string\PublishedStripe as: not yet published. See REVISION.md,
  2026-09-22}
\fi

\newcommand{\sidestripe}[1]{%
  \AddToShipoutPictureBG*{%
    \AtPageLowerLeft{%
      \put(\LenToUnit{\SideStripeX},\LenToUnit{\SideStripeY}){%
        \rotatebox{90}{%
          \color[gray]{0.55}%
          \fontsize{7}{8}\selectfont\ttfamily
          \ifSeriesDraft\SeriesDraftWord\SideStripeSep\fi%
          \SideStripeRepo\SideStripeSep results/#1\SideStripeSep
          \PublishedStripe\SideStripeSep last updated \BuildDateLong}}}}}

% CHANGE OF 2026-09-22, on the owner's word, to the one file all the papers of the series share.
%
% The stripe now carries TWO dates and names the folder as results/ rather than research/:
%
%   DRAFT . evilaliv3/vesuvius-challenge-pipeline . results/<folder> . <published> . last updated <date>
%
% \PublishedStripe comes from a published-date.tex beside the paper and is the day the folder
% entered the PUBLIC repository, written by hand once and never taken from the clock; a folder
% that is not published yet defines it as «not yet published». \BuildDateLong is the other
% quantity and keeps coming from SOURCE_DATE_EPOCH. A build refuses a publication date that is
% after the build date.
%
% EVERY PAPER THAT INPUTS THIS FILE MUST NOW ALSO INPUT A published-date.tex, or \PublishedStripe
% is undefined and the build stops. The three published papers of paper/00001 to 00003 are NOT
% rebuilt today: they are recompiled with their true publication dates at their next revision,
% which is the owner's decision of 2026-09-22 and is written in REVISION.md with its date.
, and the build of each one refuses a paper that does not, so none of them can
% miss the switch and no paper of the series can be marked while another is not.
%
% Off, this file draws no watermark and loads no package it did not load before, and each paper
% rebuilds to the bytes it had before the mark existed. That was checked on 2026-09-18 by building
% all three with the switch off and comparing the sha256 with the published files, not assumed.
\newif\ifSeriesDraft
\SeriesDraftfalse

\usepackage{eso-pic}
\usepackage{xcolor}

% The build writes build-date.tex and every paper of the series inputs it before this file. If it
% is missing, the date is not declared anywhere, and the build stops here rather than reaching a
% page that quietly carries no date or, worse, one taken from the clock.
\ifdefined\BuildDateLong\else
  \PackageError{sidestripe}{The command BuildDateLong is not defined: input build-date first}
    {The build writes build-date.tex from SOURCE_DATE_EPOCH. Run the paper's build script.}
\fi

\newlength{\SideStripeX}\setlength{\SideStripeX}{0.42in}
\newlength{\SideStripeY}\setlength{\SideStripeY}{1.05in}
% The separator is a middle dot: the house rule forbids a dash as punctuation, and a comma would
% read as a list rather than as the three fields of an identifier.
\newcommand{\SideStripeSep}{\,\textperiodcentered\,}
\newcommand{\SideStripeRepo}{evilaliv3/vesuvius-challenge-pipeline}

% The watermark, drawn only while the switch is on.
%
% Why a mark on every page and not a banner on the first. A reader almost never meets these papers
% whole: a page is printed, a screenshot is pasted into a message, a section is forwarded. A mark
% that lives on the first page alone is gone at that moment, and the page that travels is the one
% that will be quoted back. A running head would reach every page too, but it belongs to the
% galley: in this two column class it would move the text block, change where the columns break
% and cost a page, and it would have to be removed from the galley again to get the final paper
% back. A shipout picture costs nothing of the kind.
%
% Why the FOREGROUND and not the background. It was drawn in the background until 2026-09-18, and
% there it was invisible wherever a figure sat, because every figure in these papers fills its own
% box with opaque white before it draws anything. A figure is exactly the part of a page somebody
% screenshots, so the mark was missing from the part most likely to travel alone. On top of
% everything it reaches the whole page, figures included.
%
% What keeps it from spoiling the page. It is drawn at \SeriesDraftOpacity of full opacity, so what
% lands on the paper is a tint and not a shape: over white it is the lightest grey a PDF can carry,
% and over black text it changes nothing at all, because a tint of black over black is black. The
% text underneath is therefore not dimmed in any way a reader can measure. Nothing is added to the
% galley either, so the page count, the line breaks and the column breaks are exactly those of the
% final paper, and the gate on overfull boxes sees nothing new.
%
% The colour is black and not a grey, because at this opacity the colour is the only thing left to
% spend: the tint over white is 1 minus the opacity of the way from white to the ink, so black is
% the darkest mark two per cent can buy. A light grey at two per cent would be indistinguishable
% from the paper.
%
% \resizebox scales a box set at the normal size instead of asking for a font at 80pt, which the
% Computer Modern shapes of this class do not have: LaTeX would substitute the nearest design size
% and print a mark several times smaller than the one asked for, with a font warning the gates do
% not read. Scaling the box cannot fail that way, and it is a linear transform in the PDF, so the
% result is the same bytes on every build.
\ifSeriesDraft
  % transparent.sty writes the constant alpha into an ExtGState, which is what pdfTeX needs to put
  % ink on top of ink without hiding it. It is loaded only here, so a paper built with the switch
  % off carries no transparency machinery and no extra resource at all.
  \usepackage{transparent}
  \newcommand{\SeriesDraftWord}{DRAFT}
  \newcommand{\SeriesDraftOpacity}{0.10}
  \AddToShipoutPictureFG{%
    \AtPageCenter{%
      \makebox(0,0){%
        \rotatebox{45}{%
          \transparent{\SeriesDraftOpacity}%
          \resizebox{0.68\paperwidth}{!}{%
            \color{black}\sffamily\bfseries\SeriesDraftWord}}}}}
\fi

% THE CASE OF THE SERIES: the headings and the table captions read as the rest of the article.
%
% IEEEtran sets two things in small capitals, and a reader sees them as ALL CAPITALS shouted in the
% middle of a page whose every other line is ordinary sentence case:
%
%   the table captions, IEEEtran.cls line 2778, which puts the caption text through \scshape;
%   the section headings, IEEEtran.cls line 5473, whose heading font ends in \scshape.
%
% Those are the two the responsible person asked for on 2026-09-18, and they are the only two that
% are live in this class option. The other \scshape sites in the class, lines 2737, 5468, 5501 and
% 5790, belong to the conference and computer society options, which these papers do not use, and
% the figure captions never had it: figure captions were already sentence case while table captions
% were not, so this change also makes the two agree, which they did not before.
%
% What is deliberately NOT changed: the drop cap at the start of each section, IEEEtran.cls lines
% 5835 and 5841, which sets the large initial and the rest of its first word through
% \MakeUppercase. A large initial followed by a word in capitals is a typographic convention and
% not shouting, and it is a bigger aesthetic move than was asked for. One line in this file takes
% it off if it is ever wanted: \renewcommand{\IEEEPARstartWORDCAPSTYLE}{\relax}.
%
% The class file is NOT edited. It is a third party class, vendored in three trees, and a modified
% copy under the same name is what makes a build unreproducible for somebody else. etoolbox patches
% the two macros from the preamble instead, so one change here reaches all three papers exactly as
% the draft switch does.
%
% Both patches fail the build if they do not apply. That matters more than it looks: if the class
% is ever updated and the macro no longer matches, a silent failure would bring the capitals back
% and nobody would know which build changed them.
\usepackage{etoolbox}
\makeatletter
\patchcmd{\@makecaption}{\scshape}{}{}{%
  \PackageError{sidestripe}{Cannot take the small capitals out of the table captions}
    {IEEEtran's \string\@makecaption no longer contains \string\scshape where it did at line
     2778. Read the class and fix the patch rather than removing it.}}
\patchcmd{\section}{\scshape}{}{}{%
  \PackageError{sidestripe}{Cannot take the small capitals out of the section headings}
    {IEEEtran's \string\section no longer contains \string\scshape where it did at line 5473.
     Read the class and fix the patch rather than removing it.}}
\makeatother

% Mechanics of the stripe. eso-pic's \AddToShipoutPictureBG* draws in the background layer of the
% current page only, which is what keeps the body untouched: the text block is not moved by a
% single point, because nothing is added to the galley. \rotatebox comes from graphicx, which the
% papers already load for their figures, and the grey is xcolor. tikz with `remember picture,
% overlay' would do the same thing at the price of a second compilation pass and a much larger
% dependency, and a raw \rotatebox in a zero size box inside the body would have to be placed by
% hand relative to the text and would move with it.
%
% Geometry. The stripe sits \SideStripeX from the left edge of the paper and starts \SideStripeY
% from the bottom, reading upwards. Those are far enough inside the paper that a printer's trim
% does not eat them and far enough from the text block that they cannot collide with it; both are
% named here so that the three papers cannot drift apart.
% If the paper that inputs this file has no published-date.tex, say so in words rather than
% letting TeX report an undefined control sequence from inside a rotated box. The three papers of
% paper/00001 to 00003 are in exactly that state on 2026-09-22 and are not rebuilt today; whoever
% builds one next is told here what to add and why, instead of debugging it.
\ifdefined\PublishedStripe\else
  \errmessage{sidestripe.tex: this paper inputs the stripe but defines no \string\PublishedStripe.
  Add a published-date.tex beside the paper and \string% The day this folder entered the PUBLIC repository, written here by hand once and never taken
% from the clock: a build that read the date from the machine would move it every time and a
% front page would claim a day it was not published on. The last updated date beside it in the
% stripe is a different quantity and comes from SOURCE_DATE_EPOCH.
%
% This folder is NOT in the public repository yet: on 2026-09-22 the only folder under
% results/ of evilaliv3/vesuvius-challenge-pipeline is e89eb7ad-umbilicus-estimator. So the
% stripe carries «not yet published» rather than a date, and when the folder is published the
% two lines below are replaced with the day it was, in this shape:
%
%   \newcommand{\PublishedDateIso}{2026-09-22}
%   \newcommand{\PublishedStripe}{published 22 September 2026}
%
% build.sh reads \PublishedDateIso and refuses a date in the future.
\newcommand{\PublishedDateIso}{}
\newcommand{\PublishedStripe}{not yet published}
. It holds the day
  the folder entered the PUBLIC repository, written by hand and never from the clock; a folder not
  yet published defines \string\PublishedStripe as: not yet published. See REVISION.md,
  2026-09-22}
\fi

\newcommand{\sidestripe}[1]{%
  \AddToShipoutPictureBG*{%
    \AtPageLowerLeft{%
      \put(\LenToUnit{\SideStripeX},\LenToUnit{\SideStripeY}){%
        \rotatebox{90}{%
          \color[gray]{0.55}%
          \fontsize{7}{8}\selectfont\ttfamily
          \ifSeriesDraft\SeriesDraftWord\SideStripeSep\fi%
          \SideStripeRepo\SideStripeSep results/#1\SideStripeSep
          \PublishedStripe\SideStripeSep last updated \BuildDateLong}}}}}

% CHANGE OF 2026-09-22, on the owner's word, to the one file all the papers of the series share.
%
% The stripe now carries TWO dates and names the folder as results/ rather than research/:
%
%   DRAFT . evilaliv3/vesuvius-challenge-pipeline . results/<folder> . <published> . last updated <date>
%
% \PublishedStripe comes from a published-date.tex beside the paper and is the day the folder
% entered the PUBLIC repository, written by hand once and never taken from the clock; a folder
% that is not published yet defines it as «not yet published». \BuildDateLong is the other
% quantity and keeps coming from SOURCE_DATE_EPOCH. A build refuses a publication date that is
% after the build date.
%
% EVERY PAPER THAT INPUTS THIS FILE MUST NOW ALSO INPUT A published-date.tex, or \PublishedStripe
% is undefined and the build stops. The three published papers of paper/00001 to 00003 are NOT
% rebuilt today: they are recompiled with their true publication dates at their next revision,
% which is the owner's decision of 2026-09-22 and is written in REVISION.md with its date.
